\chapter{Multiple Access and the Cellular Concept}
\label{ch:multipleaccess}
\providecommand{\chTwentyFigH}[3][0.92]{\begin{figure}[H]\centering\includegraphics[width=#1\linewidth]{figs/#2}\caption{#3}\label{fig:#2}\end{figure}}
\providecommand{\chTwentyPhotoH}[3][0.6]{\begin{figure}[H]\centering\includegraphics[width=#1\linewidth]{figs/photos/#2}\caption{#3\mdccredit{#2}}\label{fig:#2}\end{figure}}

\begin{chapterintro}
Walk out of a sold-out stadium with sixty thousand other people and, within a minute or two,
tens of thousands of phones wake up and ask the network for attention at once: to post a photo,
book a ride, call home. Every one of them is shouting into the same few tens of megahertz of
spectrum. Remarkably, most of them get through. This chapter is about how.

Everything so far has concerned one link: one transmitter, one receiver, one channel. Real radio
systems are crowded, and the central problem of this chapter is \emph{sharing}. We start with the
two oldest questions: how one terminal can talk and listen at the same time (duplexing), and how
many terminals can split one channel without stepping on each other (FDMA, TDMA, CDMA, OFDMA,
SDMA and the non-orthogonal schemes). Then we let terminals transmit whenever they please, which
is the territory of random access: ALOHA and its famous $1/e$, carrier sense and Ethernet, the
hidden-node problem, Wi-Fi's collision avoidance and Bianchi's model, and the random-access channel
that every phone uses to say ``I am here''.

The second half covers the idea that turned mobile telephony from a luxury for a few hundred
people per city into a service for billions: the \emph{cellular concept}. Hexagonal cells and
frequency reuse, the co-channel interference that reuse creates, Erlang's traffic theory for
deciding how many channels a cell needs, interference management, the stochastic-geometry view of
network capacity, handover and paging for users who move, and the schedulers that decide, every
millisecond, who gets each resource. A short tour of radio-access network architecture, from the
GSM base-station controller to O-RAN, hands over to the next chapter's history of the cellular
generations.
\end{chapterintro}

\mdcfig{ch20_timeline}{Seventy years of sharing the air. Each milestone in this chapter added a new
way to divide a channel among many users: frequency (MTS, AMPS), time (GSM), codes (IS-95), carrier
sense and collision avoidance (ALOHAnet, Ethernet, 802.11), the time--frequency grid (LTE, NR) and,
most recently, open software control of the whole radio network (O-RAN).}

%======================================================================
\section{Sharing a scarce resource}
\label{sec:ch20:sharing}
%======================================================================

Picture a long dinner table with twenty guests and only one rule: nobody may raise their voice.
If everyone talks at once, nobody hears anything. Yet real dinner parties work. People take turns,
split into little conversations at different ends of the table, lean towards the person they are
talking to, or switch to a language the neighbours do not share. Without thinking about it,
dinner guests reinvent every multiple-access scheme in this chapter.

A radio channel is that dinner table. Every transmission in a band raises the noise floor for
every other receiver within range, and the band cannot be made larger: the spectrum below about
6\,GHz that propagates well through walls and foliage is finite, licensed at auction for tens of
billions of dollars, and shared by everything from television to radar. A system designer facing
$K$ users and one band has to answer three questions.
\begin{enumerate}
\item \textbf{Duplexing.} How does each terminal transmit and receive? A transmitter's own
signal at its receiver input is typically more than 100\,dB stronger than the signal it is
trying to hear, so the two directions must be kept apart somehow.
\item \textbf{Multiple access.} How are the resources divided among users? The division may
be fixed (each user gets a channel for the duration of a call), scheduled (a central
controller assigns resources moment by moment), or random (users transmit when they have
data and resolve the collisions afterwards).
\item \textbf{Reuse.} How can the same resources be used again somewhere else? A transmitter's
signal fades with distance, so a frequency used in one place can be reused far enough away.
The art of deciding what ``far enough'' means is the cellular concept.
\end{enumerate}

\mdcfig{ch20_dinner}{Multiple access at the dinner table. Guests can take turns at one table (time
division), split into separate tables (frequency division), or all talk at once in different
languages, each listener tuning in to their own (code division). Radio engineers use all three, and
combinations of them.}

\begin{analogy}[A dinner-table conversation]
Taking turns is TDMA: one speaker at a time, everybody hears clearly, but each guest waits for
a slot. Sitting at different tables is FDMA: conversations run in parallel and never mix, but
each table is small and a guest stuck at a dull table cannot borrow a seat elsewhere. Speaking
different languages is CDMA: everyone talks at once, and each listener picks out their own language
from a background babble that grows as more people join. Leaning towards your partner and speaking
quietly is SDMA, separation in space. The limits of the analogy: radio users cannot see each other,
so they cannot politely notice who is speaking, and the ``languages'' of CDMA are mathematical codes
that only work if every voice arrives at nearly the same loudness.
\end{analogy}

\mdcfig[0.98]{ch20_access_dims}{The time--frequency plane divided four ways among four users.
FDMA assigns bands, TDMA assigns slots, CDMA lets every user occupy the whole plane with a different
code (drawn as stacked layers), and OFDMA assigns blocks of the grid that a scheduler can change
from one slot to the next.}

\subsection{Orthogonality: one idea, many dimensions}

Every classical multiple-access scheme rests on one idea from Chapter~\ref{ch:modulation}: signals
that are \emph{orthogonal} do not interfere. In words: if two users' waveforms have zero overlap
when you multiply them together and integrate, a receiver tuned to one of them sees nothing of the
other. Formally, suppose user $k$ sends $x_k(t)=\sum_m a_{k,m}\,\phi_{k,m}(t)$ using a set of basis
waveforms $\phi_{k,m}$. If the waveforms of different users are orthogonal,
\begin{equation}
\int\phi_{k,m}(t)\,\phi^*_{\ell,n}(t)\,dt=0\qquad(k\ne\ell),
\label{eq:ch20:orth}
\end{equation}
then a receiver that correlates against user $k$'s waveforms sees none of user $\ell$'s energy,
and each user behaves as if alone. The dimension in which orthogonality is arranged names the
scheme (Figure~\ref{fig:ch20_access_dims}).
\begin{itemize}
\item \term{Frequency-division multiple access} (FDMA) gives each user a different frequency
band. Waveforms with disjoint spectra are orthogonal.
\item \term{Time-division multiple access} (TDMA) gives each user different time slots on a
common carrier. Waveforms with disjoint supports in time are orthogonal.
\item \term{Code-division multiple access} (CDMA) lets all users share time and frequency but
gives each a different spreading code, orthogonal or nearly so (Chapter~\ref{ch:spreadspectrum}).
\item \term{Orthogonal frequency-division multiple access} (OFDMA) divides a time--frequency
grid of OFDM resource elements among users (Chapter~\ref{ch:ofdm}); it is FDMA and TDMA at once,
with the subcarriers packed as tightly as orthogonality allows.
\item \term{Space-division multiple access} (SDMA) separates users by their spatial signatures
at an antenna array (Chapter~\ref{ch:mimo}).
\end{itemize}

Here is a fact that surprises many students. The time--bandwidth product of a channel of bandwidth
$W$ and duration $T$ is about $2WT$ real dimensions (Chapter~\ref{ch:signals}). FDMA, TDMA and
orthogonal CDMA are simply three different bases for the same signal space, and in an ideal channel
they support exactly the same total rate: $K$ users sharing $2WT$ dimensions get $2WT/K$ each,
however the dimensions are sliced. Cutting a cake into horizontal layers or vertical wedges does
not change how much cake there is. The differences between the schemes lie entirely in how they
behave when the ideal assumptions fail: when the channel has multipath that smears time, Doppler
that smears frequency, or users whose signals arrive with very different powers, and when traffic
is bursty rather than continuous.

\begin{keyidea}[The three ways to lose orthogonality]
In practice orthogonality is always approximate, and every real system can be understood by
asking how it fails. FDMA loses it through imperfect filters and oscillator drift, so it needs
guard bands. TDMA loses it through timing errors and propagation delay, so it needs guard times
and timing advance. CDMA loses it through multipath and asynchronous arrival, which destroy the
code orthogonality, so it needs power control. OFDMA loses it through Doppler, frequency offsets
and late echoes, so it needs a cyclic prefix and tight synchronization. The overhead each
scheme pays---guard bands, guard times, pilots, prefixes, power-control signalling---is the
price of approximate orthogonality.
\end{keyidea}

\mdcfig{ch20_one_tower}{The cellular idea in one picture. (a) Before cells, one powerful transmitter
covered a city and each channel carried one call for the whole metropolitan area. (b) With small cells
and a reuse pattern of three channel sets (A, B, C), the same twelve channels serve 76 simultaneous
calls in 19 cells, and capacity keeps growing as cells shrink.}

\subsection{Before cells: one tower for a whole city}

The first mobile telephones show what happens when you divide a band but never reuse it.

\mdcpairany[htbp]{ch20_car_phone}{A General Electric car radiotelephone control head with a rotary
dial, from the era of dial-up mobile service, now in a telephone museum. The transceiver itself
lived in the car's boot; the lamps show when the channel is in use (``busy'') and when the
radio is transmitting.}{ch20_mta_phone}{A Swedish car telephone for the MTA system of 1956
(Tekniska museet, Stockholm): racks of valves and relays behind a dial and handset. Every
early mobile system was FDMA with a handful of channels per city.}

\begin{history}[Mobile telephony before cells]
The first public mobile telephone service in the United States opened in St.~Louis on 17~June
1946, when Southwestern Bell and AT\&T connected a handful of car-mounted radiotelephones near
150\,MHz to the telephone network. This Mobile Telephone Service (MTS) was a single high-power
transmitter on a tall tower covering a whole metropolitan area, with a few channels, a
push-to-talk handset and an operator who placed every call. It was an FDMA system in its
purest form, and the channels were not reused anywhere in the city. The Improved Mobile
Telephone Service (IMTS) of 1964 added automatic channel selection, direct dialling and
full-duplex operation, and later added channels near 450\,MHz, but kept the
one-big-transmitter architecture. The result was a service for the rich and the patient. A
figure widely quoted in textbooks (for example by Rappaport) is that in 1976 the Bell system's
mobile service in New York City had only twelve channels, served about 550 customers and had a
waiting list of more than 3700; call blocking in the busy hour was severe. More channels could
not fix the problem, because every channel was used once in the whole city. The fix, the
cellular concept of Section~\ref{sec:ch20:cellular}, had been described on paper in 1947.
\end{history}

Figure~\ref{fig:ch20_one_tower} puts the 1976 New York numbers in perspective. Twelve channels on
one tower carry at most twelve calls, whether the city has a million people or ten million. Cut the
same city into small cells, give each cell a third of the channels, and every cluster of three cells
reuses all twelve. Nineteen cells already carry 76 calls; a thousand cells would carry four thousand.
Nothing about the radio has improved; only the geometry has changed. Hold on to that picture: it is
the whole of Section~\ref{sec:ch20:cellular} in one drawing.

%======================================================================
\section{Duplexing: talking and listening}
\label{sec:ch20:duplex}
%======================================================================

Try listening for a whisper from across the room while you are shouting. That is the duplexing
problem, and the numbers are brutal. A handset transmitting at $+23$\,dBm may need to receive a
signal at $-100$\,dBm (Figure~\ref{fig:ch20_isolation}a). The 123\,dB difference is a power ratio of
two million million, more than the isolation of any practical antenna arrangement, so a terminal
cannot simply transmit and receive on the same frequency at the same time. Two solutions have
dominated for half a century, and a third is emerging.

\mdcfig{ch20_isolation}{Why duplexing is hard. (a) A handset's own transmitter arrives at its antenna
port about 123\,dB stronger than the signal it is trying to hear, which sits just above the thermal
noise. (b) In-band full duplex must close that gap with roughly 110\,dB of self-interference
suppression, stacked from antenna isolation, an analog canceller and a digital canceller; the split
shown is one illustrative budget, and real designs divide it differently.}

\subsection{Frequency-division duplexing}

\term{Frequency-division duplexing} (FDD) is a two-lane road: uplink traffic in one lane, downlink in
the other, a painted gap between them. It uses two paired bands, one for the uplink (terminal to
base station) and one for the downlink, separated by a fixed \term{duplex spacing}. A
\term{duplexer}, a pair of sharp band-pass filters joined at the antenna, steers the transmit
signal to the antenna and keeps it out of the receiver (Chapter~\ref{ch:sdr}). AMPS used a
45\,MHz spacing (uplink 824--849\,MHz, downlink 869--894\,MHz); so do GSM-900 and LTE band~5,
and most FDD bands below 3\,GHz have spacings between 30 and 400\,MHz. The duplexer has to
attenuate the transmitter's own signal and its noise in the receive band by roughly 50\,dB or
more, and the gap between the bands (the \term{duplex gap}) must be wide enough for affordable
filters to roll off: a narrow gap is a filter problem, a wide one is wasted spectrum.

FDD's strengths are continuity and simplicity. Both directions are always available, so there
is no switching, no guard time, and no need for neighbouring base stations to coordinate their
timing. Its weakness is rigidity. Spectrum must come in paired, equal-sized blocks, while
traffic is strongly asymmetric: mobile networks typically carry several times more data on the
downlink than the uplink, so a symmetric FDD pair wastes much of its uplink half. And because
the two directions are on different frequencies, the downlink channel cannot be learned from
uplink measurements; the terminal must measure and report it.

\begin{figure}[htbp]\centering
\begin{tikzpicture}[x=1cm,y=1cm,font=\sffamily\scriptsize]
% FDD
\node[anchor=west,font=\sffamily\small\bfseries,text=navy] at (0,3.55) {FDD};
\fill[navy!70] (0.9,2.45) rectangle (6.0,2.95); \node[text=white] at (3.45,2.7) {downlink band};
\fill[accent!70] (0.9,1.55) rectangle (6.0,2.05); \node[text=white] at (3.45,1.8) {uplink band};
\draw[<->,gray] (6.2,1.8) -- node[right,align=left]{duplex\\spacing} (6.2,2.7);
\draw[->] (0.8,1.2) -- (6.3,1.2) node[right]{time};
\draw[->] (0.8,1.2) -- (0.8,3.2) node[left]{freq.};
% TDD
\node[anchor=west,font=\sffamily\small\bfseries,text=navy] at (8.0,3.55) {TDD (pattern DDDSU)};
\foreach \i in {0,1,2} {\fill[navy!70] (8.9+\i*0.9,1.9) rectangle (9.75+\i*0.9,2.6); \node[text=white] at (9.32+\i*0.9,2.25){D};}
\fill[navy!70] (11.6,1.9) rectangle (12.05,2.6);
\fill[gray!25] (12.05,1.9) rectangle (12.25,2.6);
\fill[accent!70] (12.25,1.9) rectangle (12.45,2.6);
\node at (11.83,2.8) {S};
\fill[accent!70] (12.5,1.9) rectangle (13.35,2.6); \node[text=white] at (12.92,2.25){U};
\draw[->] (8.8,1.2) -- (13.6,1.2) node[right]{time};
\draw[->] (8.8,1.2) -- (8.8,3.2) node[left]{freq.};
\draw[gray,->] (13.15,3.05) node[right]{guard period} -- (12.17,2.62);
% timing detail
\begin{scope}[yshift=-0.2cm]
\node[anchor=west,font=\sffamily\small\bfseries,text=navy] at (0,0.85) {Why the guard period: the view at the base station and at a terminal at distance $d$};
\draw (0.9,-0.4) -- (13.4,-0.4); \node[anchor=east] at (0.85,-0.4) {BS};
\draw (0.9,-1.4) -- (13.4,-1.4); \node[anchor=east] at (0.85,-1.4) {UE};
\fill[navy!70] (1.0,-0.3) rectangle (5.5,-0.5); \node[above] at (3.2,-0.3) {downlink as sent};
\fill[navy!40] (2.2,-1.3) rectangle (6.7,-1.5); \node[below] at (4.4,-1.5) {downlink as received ($+d/c$)};
\fill[accent!70] (8.0,-1.3) rectangle (12.5,-1.5); \node[below] at (10.25,-1.5) {uplink sent early by the timing advance $2d/c$};
\fill[accent!40] (9.2,-0.3) rectangle (13.3,-0.5); \node[above] at (11.2,-0.3) {uplink arrives on the BS grid};
\draw[dashed,gray] (5.5,-0.4) -- (6.7,-1.4); \draw[dashed,gray] (8.0,-1.4) -- (9.2,-0.4);
\draw[<->,thick,navy] (6.7,-2.15) -- node[below]{switch time + slack} (8.0,-2.15);
\draw[<->,thick,accent] (5.5,0.0) -- node[above]{$T_{\rm GP}\ge 2d/c+T_{\rm switch}$} (9.2,0.0);
\end{scope}
\end{tikzpicture}
\caption{Duplexing. FDD uses two paired bands continuously; TDD alternates one band in time, here
in the three-downlink, one-special, one-uplink pattern common in 5G mid-band networks. The special
slot contains downlink symbols, a guard period and uplink symbols. Bottom: the guard period must
absorb the round-trip propagation delay of the most distant terminal plus its receive-to-transmit
switching time.}
\label{fig:ch20_duplex}
\end{figure}

\subsection{Time-division duplexing}

\term{Time-division duplexing} (TDD) is a single-lane bridge with a traffic light: everyone uses the
whole width, but in turns. It uses one band for both directions, alternating in time. The ratio of
downlink to uplink time can match the traffic, unpaired spectrum can be used, and the radio needs a
fast switch instead of a duplexer. Most 5G spectrum above 2.5\,GHz is unpaired and operated in TDD,
as is Wi-Fi (in effect: a station and its access point take turns), DECT, Bluetooth and the Chinese
TD-SCDMA and TD-LTE systems.

TDD has two subtleties. The first is the \term{guard period}, the empty gap that lets the bridge
clear before traffic flows the other way. A terminal at distance $d$ hears the end of a downlink burst
$d/c$ late, and its uplink burst must arrive at the base station on time, so it must start
transmitting $d/c$ early (its \term{timing advance}). Between the end of the downlink as seen at the
base station and the start of the uplink, the terminal must therefore fit a round trip of $2d/c$ plus
the time its own radio needs to switch from receive to transmit. The guard period sets the largest cell:
\begin{equation}
T_{\mathrm{GP}}\;\ge\;\frac{2d_{\max}}{c}+T_{\mathrm{switch}}
\qquad\Longrightarrow\qquad
d_{\max}\approx\frac{c\,(T_{\mathrm{GP}}-T_{\mathrm{switch}})}{2}.
\label{eq:ch20:gp}
\end{equation}
In words: every microsecond of guard buys 150\,m of cell radius. A 100\,km rural cell would need a
guard of more than 0.67\,ms, an expensive overhead in a frame of a few milliseconds. TDD therefore
favours small cells, and long-range systems prefer FDD (Figure~\ref{fig:ch20_tdd_range}).

The second subtlety is that every base station in an area must use the same frame timing and
the same uplink--downlink split. If one base station transmits while its neighbour is
receiving, the transmitting base station, high on a mast with line of sight to the other, drowns
the weak uplink signals at its neighbour. This \term{cross-link interference} is so severe that
TDD networks are synchronised to GNSS time to within about a microsecond, and operators in
adjacent TDD blocks are usually required by regulators to adopt compatible frame structures.

\mdcfig[0.5]{ch20_tdd_range}{The worked example generalised: maximum cell range against the number
of guard symbols, with 10\,$\mu$s of switching time, for three NR numerologies. Wider subcarrier
spacing means shorter symbols, so millimetre-wave numerologies (120\,kHz) suit only small cells.}

\begin{worked}[How large a cell can a 5G TDD guard support?]
A 5G NR carrier at 3.5\,GHz uses 30\,kHz subcarrier spacing, so a slot of 14 OFDM symbols lasts
0.5\,ms and each symbol (with its cyclic prefix) about 35.7\,$\mu$s. A widely used configuration
repeats the pattern DDDSU every 2.5\,ms, with the special slot S split as 10 downlink symbols,
2 guard symbols and 2 uplink symbols. What maximum cell range does the guard allow, and what
does the pattern cost?

The guard is $2\times35.7\approx71.4\,\mu$s. Allowing about 10\,$\mu$s for the terminal's and
base station's transmit--receive switching, \eqref{eq:ch20:gp} gives
$d_{\max}\approx3\times10^8\times(61.4\times10^{-6})/2\approx9.2$\,km, ample for urban and
suburban mid-band cells, whose range is limited by the uplink link budget long before that. The
time split is $(3+10/14)/5\approx74\%$ downlink, $(1+2/14)/5\approx23\%$ uplink and 3\% guard. A
rural deployment that needs 30\,km cells would need about 4 guard symbols, at a further cost of
about 3\% of capacity.
\end{worked}

\subsection{Channel reciprocity}

TDD has one property that FDD lacks and that has become very valuable: \term{channel reciprocity}.
Electromagnetic propagation is reciprocal---the road from A to B has the same potholes as the road
from B to A---so when uplink and downlink use the same frequency within the channel's coherence time
(Chapter~\ref{ch:channels}), the propagation channel in one direction is the transpose of the channel
in the other. A base station can then estimate the downlink channel to every antenna of every terminal
from uplink pilots (5G NR's sounding reference signals) and compute its precoder without any feedback.
With 64 or more antennas and many users, feedback-based channel acquisition would be prohibitively
expensive, which is why massive MIMO (Chapter~\ref{ch:mimo}) is predominantly a TDD technology.
Reciprocity holds for the propagation channel only: the transmit and receive RF chains of each radio
differ in gain and phase and must be calibrated, usually by having the array's antennas sound one another
(Figure~\ref{fig:ch20_reciprocity}).

\begin{figure}[htbp]\centering
\begin{tikzpicture}[font=\sffamily\scriptsize,>={Stealth[length=2mm]},
  blk/.style={draw=navy,fill=navy!7,thick,rounded corners=1pt,minimum height=0.7cm,align=center}]
\node[blk] (ue) at (0,0) {terminals\\(1 or 2 antennas each)};
\node[blk] (est) at (4.6,0) {base station:\\estimate $\mathbf H_{\rm UL}$ from pilots};
\node[blk] (pre) at (9.2,0) {precoder from\\$\mathbf H_{\rm DL}=\mathbf H_{\rm UL}^{T}$};
\draw[->,thick,accent] (ue) -- node[above]{uplink pilots (SRS)} (est);
\draw[->,thick] (est) -- node[above]{reciprocity} node[below,text=gray]{+ RF calibration} (pre);
\draw[->,thick,navy] (pre.south) -- ++(0,-0.6) -| node[pos=0.25,below]{downlink beams, same band, within the coherence time} (ue.south);
\end{tikzpicture}
\caption{Channel reciprocity in TDD. Uplink pilots tell the base station the channel; because uplink and downlink
share a frequency, the downlink channel is its transpose, and the precoder can be computed with no feedback from the
terminals---the reason massive MIMO is mostly a TDD technology.}
\label{fig:ch20_reciprocity}
\end{figure}

\subsection{Dynamic TDD, half-duplex FDD and full duplex}

\term{Dynamic TDD} changes the uplink--downlink split from slot to slot to follow the traffic of
each cell. 5G NR supports this fully: the slot format can be set semi-statically or signalled
dynamically in downlink control information. In practice, dynamic TDD is limited to small or
isolated cells, because of the cross-link interference described above.

\term{Half-duplex FDD} uses paired bands but never transmits and receives at the same instant,
so the terminal needs a switch rather than a duplexer. GSM terminals have always worked this
way: a handset transmits three timeslots after it receives, so the uplink and downlink slots
never overlap and the switch can replace most of the duplexer's isolation. Half-duplex FDD is
an option for the low-cost LTE-M and NB-IoT device categories and for 5G's reduced-capability
(RedCap) devices, where removing the duplexer saves cost and size.

\term{In-band full duplex} transmits and receives on the same frequency at the same time,
potentially doubling spectral efficiency. It requires suppressing the transmitter's own signal at
its receiver to near the noise floor: about 110\,dB for a Wi-Fi-class radio, achieved by
combining antenna isolation, an analog canceller that subtracts a tapped and delayed copy of the
transmit signal before the low-noise amplifier, and a digital canceller that models the
remaining linear and nonlinear leakage (Figure~\ref{fig:ch20_isolation}b). Laboratory
demonstrations, notably by Bharadia, McMilin and Katti at Stanford in 2013, reached about this
level. It is like wearing noise-cancelling headphones that must cancel your own voice so perfectly
that you can hear a pin drop. A practical intermediate step is \term{subband full duplex} (SBFD),
studied by 3GPP for Release~18 and specified from Release~19, in which a TDD base station receives
uplink in one sub-band of the carrier while transmitting downlink in the rest. The base station,
with its larger size and budget, carries the burden of self-interference cancellation; terminals
remain half duplex.

\begin{inpractice}[Remote interference and the atmospheric duct]
TDD networks can be disturbed by base stations hundreds of kilometres away. Under certain
weather conditions, typically calm nights over flat land or sea, a temperature inversion forms an
\emph{atmospheric duct} that guides radio waves far beyond the horizon (Chapter~\ref{ch:channels}).
Downlink signals from a distant TDD network then arrive so late that they fall into the victim's
uplink period, raising its noise floor by many decibels. Large TD-LTE operators in China observed
this problem in their networks, and 3GPP Release~16 added \emph{remote interference management}:
an affected base station transmits a recognisable reference signal so the aggressors can identify
themselves and back off (for example by shortening their downlink). It is a striking illustration
that the guard period of~\eqref{eq:ch20:gp} assumes a propagation environment the atmosphere does
not always respect.
\end{inpractice}

%======================================================================
\section{Channelization: dividing the resource}
\label{sec:ch20:channelization}
%======================================================================

Once the directions are separated, the resources in each direction must be divided among the
users. This section compares the orthogonal schemes in engineering detail, the way the systems
that used them actually lived and died, and then introduces deliberately non-orthogonal access.

\mdcphoto[0.5]{ch20_intersection}{A traffic light: time-division access to a shared crossing. Each
direction gets the whole road for a slot, and the amber and all-red intervals are guard times for
vehicles still in the junction.}

\mdcfig{ch20_gsm_frame}{GSM's TDMA frame. Eight users take turns on one
200\,kHz carrier; each burst carries two 57-bit data blocks around a 26-bit training sequence that
the receiver uses to learn and equalize the multipath channel, and ends with a guard of 8.25 bit
periods.}

\subsection{FDMA}

FDMA is the oldest scheme and the most intuitive: each user gets a narrow channel for the
duration of a call, exactly as each radio station gets a frequency on the dial. First-generation
analog cellular systems were pure FDMA: AMPS used 30\,kHz FM channels, the European TACS 25\,kHz,
the Nordic NMT 25 or 12.5\,kHz. FDMA's virtues are that each channel is narrow, so the symbol period
is long and the channel is flat (no equalizer is needed), and that synchronization between users
is unnecessary. Its weaknesses are the cost of guard bands (adjacent-channel interference from a
nearby strong transmitter must be filtered out, which requires either guard bands or sharp, costly
filters), the need for one radio per channel at the base station, and its inflexibility: a user
who needs more bandwidth needs a second radio. Single-channel-per-carrier FDMA is still used in
satellite systems with many small terminals (Chapter~\ref{ch:satellite}) and in narrowband
land-mobile radio.

\subsection{TDMA}

A traffic light is the world's most familiar TDMA system (Figure~\ref{fig:ch20_intersection}): the
crossing is shared by giving each direction the whole road for a slice of time, with an all-red
``guard'' between slices so that late arrivals clear the junction. TDMA gives each user a short burst
on a shared, wider carrier, repeating in a \term{frame}. GSM is the canonical example: 200\,kHz carriers
with a gross rate of about 270.8\,kb/s, divided into frames of 4.615\,ms with eight timeslots of about
577\,$\mu$s each, so that one carrier serves eight full-rate voice users (Figure~\ref{fig:ch20_gsm_frame}).
The North American IS-136 (D-AMPS) put three TDMA users into each old 30\,kHz AMPS channel.

\mdcfig[0.5]{ch20_near_far}{The near--far problem. With path-loss exponent 3.5, a user 100\,m from the
base station arrives 35\,dB stronger than one at 1\,km if both transmit at full power. CDMA uplinks
therefore equalise received powers to within about a decibel by fast power control.}

TDMA has three structural costs.
\begin{itemize}
\item \textbf{Guard time and timing advance.} Bursts from users at different distances must not
overlap at the base station. GSM's timing advance, measured by the base station and sent to the
handset, has 64 steps of one bit period (about 3.7\,$\mu$s, or 550\,m of range each way), which is
the origin of GSM's famous 35\,km maximum cell radius.
\item \textbf{Equalization.} A wider carrier means shorter symbols, so multipath causes intersymbol
interference; GSM receivers equalize with a training sequence in the middle of every burst
(Chapter~\ref{ch:equalization}).
\item \textbf{Peak power.} A user who transmits one-eighth of the time must transmit at eight times
the average power to deliver the same energy. The resulting 217\,Hz burst envelope of GSM handsets
was famously audible in nearby loudspeakers.
\end{itemize}
In exchange, TDMA lets a single base-station radio serve many users, lets the handset measure
neighbouring cells in its idle slots (which made mobile-assisted handover practical), and gives
variable rates naturally by assigning more slots.

\subsection{CDMA}

In CDMA every user transmits across the whole band all the time, distinguished by a spreading
code---the cocktail party where everybody speaks a different language. In the downlink of IS-95 and
UMTS, orthogonal Walsh or OVSF codes keep the users of one cell orthogonal at the transmitter; in the
uplink, users arrive asynchronously and through different channels, so their codes are only
approximately orthogonal and each user's signal acts as noise to the others. The capacity of an
uplink CDMA cell is therefore set by interference, not by a fixed number of channels: it degrades
gracefully as users are added, it benefits automatically from voice activity and from the gaps between
bursts of data, and it allows a frequency reuse factor of one.

It also requires accurate, fast \term{power control}, because a nearby user received 30\,dB stronger
than a distant one would drown it: the \term{near--far problem} (Figure~\ref{fig:ch20_near_far}). At a
party, the guest shouting next to you in Portuguese makes it impossible to follow a quiet French
conversation across the room, even though you do not understand a word of Portuguese.
Chapter~\ref{ch:spreadspectrum} develops CDMA in detail, including the uplink capacity formula, and
Chapter~\ref{ch:cellular} tells the story of the ``CDMA wars'' of the 1990s.

\chTwentyFigH{ch20_ofdma_grid}{Why OFDMA suits a scheduler. Three users see different frequency-selective
channels across the same 600 subcarriers; giving each resource block to the user whose channel is
strongest there raises the average gain of the scheduled blocks by almost 4\,dB in this example,
compared with handing each user a fixed slice regardless of its channel.}

\subsection{OFDMA and SC-FDMA}

OFDMA, the multiple-access scheme of LTE, 5G NR, WiMAX and Wi-Fi~6 and~7, assigns users
\emph{resource blocks}, groups of subcarriers for a group of OFDM symbols, and lets a scheduler change
the assignment every slot (Chapter~\ref{ch:ofdm}). Think of it as a restaurant with a maître d' who
reassigns tables every few minutes according to who is hungry. It combines FDMA's flat subchannels,
which need only one-tap equalization, with TDMA's flexibility. Its decisive advantage is that the
scheduler can give each user the part of the band where that user's channel is currently strong,
exploiting \emph{multiuser diversity} (Section~\ref{sec:ch20:scheduling}). On the uplink, LTE uses a
DFT-precoded variant (SC-FDMA) whose lower peak-to-average power ratio helps handset power amplifiers;
5G NR offers both.

\subsection{SDMA and multi-user MIMO}

An antenna array at the base station can form several beams at once and serve several users on the
same time--frequency resource, provided their spatial signatures are sufficiently different. This is
\term{space-division multiple access}, today called \term{multi-user MIMO}. Early SDMA systems, such as
the ``smart antenna'' base stations of the late 1990s, used a few fixed beams; a modern 5G base
station with 64 transmit--receive chains forms precoded beams for up to 8 or more simultaneous downlink
layers, and the theory (Chapter~\ref{ch:mimo}) says the gain grows with the number of antennas. SDMA
is not an alternative to the other schemes but an extra dimension: users are first separated in time
and frequency by OFDMA and then, within each resource block, in space (Figure~\ref{fig:ch20_sdma}).

\begin{figure}[H]\centering
\begin{tikzpicture}[font=\sffamily\scriptsize]
\foreach \i in {0,...,7} {\fill[navy] (0,{-0.7+0.2*\i}) rectangle (0.12,{-0.62+0.2*\i});}
\node[below,text=navy] at (0.06,-0.8) {array};
\foreach \ang/\c/\n in {28/accent/1,0/green!50!black/2,-30/orange/3} {
  \fill[\c,opacity=0.22] (0.15,0) -- ++(\ang-7:4.6) arc[start angle=\ang-7,end angle=\ang+7,radius=4.6] -- cycle;
  \node[circle,fill=\c,inner sep=2.2pt] at (\ang:4.3) {};
  \node[right,text=\c] at (\ang:4.55) {user \n};}
\node[align=left,text=gray,anchor=west] at (5.8,0) {same time, same frequency:\\each user sits in its own beam\\and in the others' nulls};
\end{tikzpicture}
\caption{Space-division multiple access: one antenna array serves three users on the same resource block
by pointing a beam at each, shaped so that each beam's nulls fall on the other users.}
\label{fig:ch20_sdma}
\end{figure}

\subsection{Non-orthogonal multiple access}

All the schemes above aim for orthogonality. A different idea is to superimpose users deliberately
and separate them with a smarter receiver. Information theory has long said this is the right thing
to do: the capacity region of the Gaussian broadcast channel (one transmitter, several receivers)
is achieved by \term{superposition coding}, not by orthogonal division (Chapter~\ref{ch:infotheory}).

\begin{analogy}[A loud announcement and a quiet aside]
A teacher addresses the whole room loudly and, under the same breath, adds a quiet aside for the
pupil in the front row. The pupils at the back hear only the loud message and treat the aside as
murmur. The front-row pupil hears both: she first takes in the loud message, mentally ``subtracts''
it, and then catches the aside cleanly. That is power-domain NOMA with successive interference
cancellation. The analogy breaks where real receivers make mistakes: if the front-row pupil mishears
the loud message, her subtraction leaves garbage and the aside is lost too.
\end{analogy}

\begin{figure}[htbp]\centering
\begin{tikzpicture}[font=\sffamily\scriptsize,>={Stealth[length=2mm]},
  blk/.style={draw=navy,fill=navy!7,thick,rounded corners=1pt,minimum height=0.7cm,align=center}]
\node[blk] (rx) at (0,0) {received\\$y$};
\node[blk] (d2) at (2.6,0) {decode $s_2$\\(treat $s_1$ as noise)};
\node[blk] (re) at (5.4,0) {re-encode\\and rebuild $\hat s_2$};
\node[draw=accent,circle,thick,inner sep=1pt] (sum) at (7.4,0) {$-$};
\node[blk] (d1) at (9.6,0) {decode $s_1$\\(interference-free)};
\draw[->,thick] (rx) -- (d2); \draw[->,thick] (d2) -- (re); \draw[->,thick] (re) -- (sum);
\draw[->,thick] (sum) -- (d1);
\draw[->,thick,gray] (rx.south) -- ++(0,-0.75) -| (sum.south);
\node[text=gray] at (3.8,-0.95) {delayed copy of $y$};
\node[text=accent,anchor=west] at (10.85,0) {$\hat s_1$};
\node[text=navy] at (2.6,0.75) {far user stops here};
\end{tikzpicture}
\caption{Successive interference cancellation at the near user. The strong (far-user) layer is decoded
first, rebuilt and subtracted; the near user's own layer is then decoded as if alone. The far user runs
only the first block.}
\label{fig:ch20_sic}
\end{figure}

Consider the downlink to two users with channel gains $|h_1|^2>|h_2|^2$, so that user~1 is ``near''
and user~2 ``far''. Let $\rho_k=P|h_k|^2/N_0$ be each user's SNR if it received the whole power $P$.
In \term{power-domain NOMA} the base station sends the sum of the two users' signals, giving a
fraction $a$ of the power to the near user and $1-a$ to the far one:
\begin{equation}
x=\sqrt{aP}\,s_1+\sqrt{(1-a)P}\,s_2,\qquad 0\le a\le 1 .
\end{equation}
The far user decodes $s_2$ treating $s_1$ as noise. The near user, whose channel is better, can
decode $s_2$ too; it does so first, subtracts it (\term{successive interference cancellation},
SIC; Figure~\ref{fig:ch20_sic}), and then decodes $s_1$ with no interference from $s_2$. The achievable rates are
\begin{equation}
R_1=\log_2(1+a\rho_1),\qquad
R_2=\log_2\!\left(1+\frac{(1-a)\rho_2}{a\rho_2+1}\right).
\label{eq:ch20:noma}
\end{equation}

Sweeping $a$ from 0 to 1 traces the boundary of the rate region, which is the capacity region of
the two-user Gaussian broadcast channel. Orthogonal access, giving a fraction $t$ of the time to
user~1, achieves only the straight line $R_1=t\log_2(1+\rho_1)$, $R_2=(1-t)\log_2(1+\rho_2)$; FDMA with
optimal power allocation does somewhat better, since a narrow band concentrates power; but neither
reaches the NOMA boundary when the users' SNRs differ greatly (Figure~\ref{fig:ch20_noma}). The gain
is largest exactly when one user is near and the other far, because the far user's signal is
cheap to send to the near user (who can decode it anyway) and a small slice of power gives the near
user a high rate.

\mdcfig{ch20_noma}{Non-orthogonal multiple access. (a) Rate regions for a near user at 20\,dB and a
far user at 0\,dB SNR. Superposition coding with SIC (NOMA) encloses both time sharing and
power-optimised FDMA; at a far-user rate of 0.5\,b/s/Hz the near user gets about 5.4\,b/s/Hz with
NOMA, 4.0 with FDMA and 3.3 with time sharing. (b) The transmitted constellation when the far user's QPSK
gets 80\% of the power and the near user's QPSK 20\%: a 16-point constellation that the far user
reads as noisy QPSK and the near user decodes in two stages.}

NOMA has practical costs: SIC errors propagate, the scheme needs accurate channel knowledge to
order the users and split the power, and its gains overlap with those of multi-user MIMO and
good scheduling, which can also exploit differences between users. LTE Release~13--14 standardised
a form of it as \emph{multi-user superposition transmission} (MUST), and the ATSC~3.0 broadcast standard
uses the same principle as \emph{layered division multiplexing}, superimposing a robust mobile service
on a high-rate fixed one. The uplink analogue is simpler and much older: a base station that
decodes several users' simultaneous signals with SIC is achieving points of the multiple-access
channel's capacity region, and grant-free uplink schemes for massive IoT build on it.

\subsection{Comparing the schemes}

Table~\ref{tab:ch20:schemes} summarises the comparison. Three lessons stand out.

\begin{table}[htbp]\centering\small
\caption{Multiple-access schemes compared. ``Reuse'' is the frequency reuse factor typical of cellular
deployments of each scheme.}
\label{tab:ch20:schemes}
\begin{tabularx}{\linewidth}{@{}l X X l X@{}}
\toprule
Scheme & Orthogonality from & Main overhead & Reuse & Representative systems\\
\midrule
FDMA & disjoint bands & guard bands, one radio per channel & 7 (21 with sectors) & AMPS, NMT, TACS, SCPC satellite\\
TDMA & disjoint time slots & guard time, timing advance, equalizer training & 3--4 (with hopping) & GSM, IS-136, DECT, Iridium\\
CDMA & spreading codes & power control, soft-handover overhead & 1 & IS-95, UMTS, CDMA2000, GNSS\\
OFDMA & subcarriers + slots & cyclic prefix, pilots, tight sync & 1 (with ICIC) & LTE, NR, WiMAX, Wi-Fi 6/7\\
SDMA/MU-MIMO & spatial signatures & channel acquisition (pilots, feedback) & 1 & LTE-A, NR massive MIMO, Wi-Fi 5--7\\
NOMA & power level + SIC & receiver complexity, error propagation & 1 & LTE MUST, ATSC 3.0 LDM\\
\bottomrule
\end{tabularx}
\end{table}

First, in a single isolated cell with continuous traffic, the orthogonal schemes are equivalent in
capacity. Their practical differences come from implementation: guard overhead, equalizer
complexity, power-amplifier efficiency, and how gracefully each copes with variable rates.

Second, in a multicell network the schemes differ greatly in how they handle \emph{interference from
other cells}, and this, not single-cell efficiency, is what decided the competition between them.
Narrowband FDMA and TDMA systems need a reuse factor of 3 to 7 to keep co-channel interference
tolerable (Section~\ref{sec:ch20:cellular}), so each cell uses only a fraction of the spectrum. CDMA's
processing gain and coding allowed reuse 1, and its capacity advantage over the 2G TDMA systems came
largely from that, together with voice-activity gating. OFDMA systems also run at reuse 1 and recover
from the resulting cell-edge interference with coding, link adaptation, scheduling and interference
coordination.

Third, the best schemes are \emph{adaptive}. The change from circuit-switched voice to packet data
transformed the problem from ``divide the spectrum into fixed channels'' into ``decide, every
millisecond, who transmits on which resource''. OFDMA won because it gives a scheduler the finest
control over the time--frequency grid; Section~\ref{sec:ch20:scheduling} studies how that control is
used.

\begin{tryit}[Air interfaces side by side]
Open Lab 11 (\lab{lab11\_air\_interfaces.py}) and choose \emph{Air interfaces side by side}. Set
\emph{System} to GSM (2G) and \emph{Compare with} to LTE 20\,MHz (4G), then raise \emph{Active users}
from 1 to 12 and watch how each system carves its time--frequency plane among them: slots on a narrow
carrier for GSM, a scheduled grid for LTE. Switch \emph{Bar chart} to \emph{Spectral efficiency} and
\emph{Symbol duration} to see the trade-offs of Table~\ref{tab:ch20:schemes} as numbers.
\end{tryit}

%======================================================================
\section{Random access}
\label{sec:ch20:random}
%======================================================================

Fixed and scheduled division work well when each user's demand is steady or predictable. Computer
and sensor traffic is neither. A terminal may be silent for minutes and then have a few hundred bytes
to send immediately; reserving a channel for it wastes the channel, and asking a scheduler for
permission costs a round trip that may be longer than the transmission itself. \term{Random access}
takes the opposite approach: transmit when you have something to send, and deal with collisions when
they happen. It sounds reckless. It runs Wi-Fi, it got Ethernet started, and every phone on Earth uses
it to make first contact with its network. This section develops the theory, from ALOHA to Wi-Fi.

\chTwentyPhotoH[0.42]{ch20_alohanet}{The IEEE Milestone plaque at the University of Hawaii at M\=anoa,
dedicated in June 2021, fifty years after ALOHAnet began carrying traffic between the islands. Its
text makes the claim of this section in one sentence: channels can be shared on a large scale by
simple random-access protocols.}

\begin{history}[ALOHAnet and the birth of packet radio]
In 1968 Norman Abramson, a professor at the University of Hawaii, began building a radio network to
connect terminals on several islands to the university's central computer in Honolulu. The ALOHA
system, operational from June 1971, used two UHF channels near 407 and 413\,MHz at 9600\,b/s: one
for the hub (``Menehune'') to broadcast to all terminals, and one shared by every terminal for
sending to the hub. Abramson's protocol for the shared channel was disarmingly simple. A terminal
with a packet sent it at once. If the hub received it correctly, it broadcast an acknowledgement; if
no acknowledgement came, the terminal assumed a collision and retransmitted after a random delay.
Abramson's 1970 paper showed that this simplest of schemes could use at most $1/(2e)$, about 18\%, of
the channel; Lawrence Roberts showed in 1972 that dividing time into slots doubled it to $1/e$.
ALOHAnet linked to the ARPANET via satellite in 1972. Its ideas led directly to Ethernet, to the
random-access channels of every cellular system and satellite network, and to the tag-reading
protocols of RFID. Few papers have been as directly influential as Abramson's.
\end{history}

\begin{analogy}[Blurting out in a meeting]
Imagine a meeting with no chair. Anyone with something to say just says it. Most of the time only one
person is talking and all is well. Occasionally two people start at once; both notice that nobody
understood, both stop, and each waits a random moment before trying again. If they waited a fixed
time they would collide again forever---the randomness is what breaks the tie. That is ALOHA. The
analogy's limit: people in a meeting can hear each other start talking and politely stop, whereas
ALOHA terminals are deaf to each other and learn of a collision only from a missing acknowledgement.
\end{analogy}

\mdcfig{ch20_vulnerable}{The vulnerable period. (a) In pure ALOHA a packet is destroyed by any other
packet that starts less than one packet time before or after it: a window of $2T$. (b) In slotted ALOHA
packets start only on slot boundaries, so only a packet in the same slot can collide: the window
halves, and so does the collision rate at a given load.}

\subsection{Pure ALOHA}

Model the aggregate traffic offered to the channel---new packets and retransmissions together---as a
Poisson process of rate $G$ packets per packet time $T$. This is an approximation (retransmissions
depend on earlier collisions), but it is accurate when there are many users each contributing a small
share, and it gives the right answers. A packet that starts at time $t$ is received correctly if no
other packet starts in the interval $(t-T,t+T)$: a packet that started up to $T$ earlier would still
be on the air, and one starting up to $T$ later would overlap its tail (Figure~\ref{fig:ch20_vulnerable}).
This \term{vulnerable period} is $2T$, and the number of Poisson arrivals in it has mean $2G$, so
\begin{equation}
\Pr\{\text{success}\}=e^{-2G},\qquad S=G\,e^{-2G},
\label{eq:ch20:pure}
\end{equation}
where the \term{throughput} $S$ is the rate of successful packets, in packets per packet time, which is
also the fraction of time the channel carries useful traffic. Setting $dS/dG=e^{-2G}(1-2G)=0$ gives the
optimum $G=1/2$ and the famous maximum
\begin{equation}
S_{\max}=\frac{1}{2e}\approx0.184 .
\end{equation}
Read that number slowly: a channel with nobody in charge, at its very best, is useful less than a
fifth of the time.

\subsection{Slotted ALOHA}

If time is divided into slots of length $T$ and every transmission must begin at a slot boundary, two
packets either overlap completely or not at all. The vulnerable period shrinks to one slot and
\begin{equation}
S=G\,e^{-G},\qquad S_{\max}=\frac1e\approx0.368\quad\text{at } G=1.
\label{eq:ch20:slotted}
\end{equation}
At the optimum, a fraction $e^{-1}$ of slots are idle, $e^{-1}$ carry a success, and
$1-2e^{-1}\approx26\%$ are wasted by collisions (Figure~\ref{fig:ch20_slot_fractions}). The cost of the
doubling is a common clock, which in a satellite or cellular system is free: the terminals already
listen to a downlink that carries timing. Figure~\ref{fig:ch20_meeting} shows the difference on the
same random traffic: snapping each packet to the next slot boundary saves several packets that pure
ALOHA loses to partial overlaps.

\mdcfig{ch20_meeting}{The same 25 packet attempts, offered at $G=0.5$, under the two ALOHA rules (each
row is a station; green packets survive, red ones collide). Under pure ALOHA, partial overlaps destroy
most of the clusters; when every transmission is aligned to a slot, only packets that land in the same
slot collide. The counts in this short sample are close to the theoretical success probabilities.}

\mdcfig[0.5]{ch20_slot_fractions}{Where the slots go in slotted ALOHA. Below $G=1$ most slots are idle;
above it most are collisions. At the optimum the split is 37\% idle, 37\% success and 26\% collision.}

Figure~\ref{fig:ch20_aloha}a plots both throughput curves. Each has two regions: at light load,
throughput equals offered load and almost every packet succeeds; beyond the peak, adding traffic
\emph{reduces} the throughput, because almost every slot is a collision. That second region is not
merely inefficient; it is unstable.

\begin{tryit}[Watch packets collide]
In Lab 26 (\lab{lab26\_cellular\_system.py}) open \emph{ALOHA and CSMA, live}. With \emph{Protocol} set
to Pure ALOHA, drag \emph{Offered load G} from 0.1 to 0.5 and then to 3: watch the green packets on the
station timeline turn red and the operating point slide over the top of the throughput curve. Switch to
Slotted ALOHA at $G=1$ and count roughly how many slots are idle, successful and collided.
\end{tryit}

\subsection{The instability of ALOHA}

Why unstable? Think of a traffic jam on a ring road. Below a critical density, cars flow; above it, every
slowdown causes more slowdowns, and the jam feeds itself long after the extra cars have gone. To see the
same thing in ALOHA, abandon the Poisson assumption and track the number of \emph{backlogged} users,
those holding a packet that has collided and awaits retransmission. Following Bertsekas and Gallager,
consider $m$ users. An unbacklogged user generates a new packet in a slot with probability $q_a$ and
sends it immediately; a backlogged user retransmits in each slot with probability $q_r$. With $n$ users
backlogged, the expected number of attempts in a slot is
\begin{equation}
G(n)=(m-n)\,q_a+n\,q_r,
\end{equation}
and the probability of a success is approximately $G(n)e^{-G(n)}$. The backlog tends to grow when new
arrivals $(m-n)q_a$ exceed successes and to shrink otherwise; the difference is the \term{drift}.
Figure~\ref{fig:ch20_aloha_stability}a plots both curves against $n$. They cross three times. The left
crossing is the desired operating point: a small backlog, almost all packets succeeding. The middle
crossing is unstable: a random burst of collisions that pushes the backlog past it makes collisions more
likely still, and the system slides to the right crossing, where nearly everyone is backlogged and the
throughput is close to zero. Because the backlog performs a random walk, it will cross the unstable point
eventually, however rarely; the network then stays congested for a long time.

\mdcfig{ch20_aloha_stability}{Stability of slotted ALOHA. (a) For $m=100$ users with new-packet
probability $q_a=0.0035$ and retransmission probability $q_r=0.06$, the arrival and departure curves
cross three times: a desired stable point, an unstable point and a congested stable point. (b) Simulated
backlog with Poisson arrivals at $\lambda=0.3$ packets per slot. With a fixed retransmission probability
the backlog eventually escapes and grows without bound; Rivest's pseudo-Bayesian rule keeps it small.}

With an infinite population of users and Poisson arrivals of rate $\lambda$, matters are worse: for any
fixed $q_r$, $G(n)=\lambda+nq_r$ grows without bound, successes vanish, and slotted ALOHA is unstable for
every $\lambda>0$ (Figure~\ref{fig:ch20_aloha_stability}b, red). The cure is to make the retransmission
probability adapt to the backlog so that $G(n)$ stays near 1. If every user knew $n$, choosing
$q_r=1/n$ would do it: with $n$ people waiting, each should try with probability $1/n$ so that, on
average, exactly one speaks. Users do not know $n$, but they can estimate it from what they hear.

\begin{deeper}[Rivest's pseudo-Bayesian algorithm]
Rivest's \term{pseudo-Bayesian algorithm} (1987) updates an estimate $\hat n$ of the backlog after each
slot, using only the channel feedback (idle, success or collision) that every user hears:
\begin{equation}
\hat n\leftarrow
\begin{cases}
\max(\lambda,\ \hat n+\lambda-1) & \text{idle or success},\\
\hat n+\lambda+(e-2)^{-1} & \text{collision},
\end{cases}
\qquad q_r=\min(1,1/\hat n).
\end{equation}
The logic: each slot adds about $\lambda$ new backlogged users; an idle slot or a success suggests the
estimate was too high (or that one user left), and a collision suggests it was too low by about
$(e-2)^{-1}\approx1.39$, the expected excess given a collision when $G=1$. The rule is stable for all
$\lambda<1/e$ (Figure~\ref{fig:ch20_aloha_stability}b, blue).
\end{deeper}

Exponential backoff, which every practical system uses, is a cruder form of the same idea: each user
infers congestion from its own collisions and lengthens its backoff. Remarkably, Aldous proved in 1987
that binary exponential backoff is itself unstable in the infinite-population model, although with a
finite number of users it works well in practice.

\begin{tryit}[Make the backlog run away]
Lab 26's \emph{Slotted ALOHA's instability} experiment runs the model above live. Leave
\emph{New packets per slot $\lambda$} at 0.30, below $1/e$, and set \emph{Fixed retransmission
probability} to 0.15. Watch the fixed-$q$ backlog hover for a while and then, suddenly, run away,
while the pseudo-Bayesian trace stays near zero. Press \emph{Empty the backlogs} and try again with
$q=0.03$: the escape takes longer, but it still comes.
\end{tryit}

\chTwentyFigH[0.62]{ch20_sensor_op}{The sensor network of the worked example on the slotted-ALOHA curve.
The required throughput of 0.333 is met at two offered loads: the stable point $G=0.62$, where each report
needs about 1.9 attempts, and a congested point $G\approx1.5$ that a burst of collisions can tip the network into.}

\begin{worked}[Slotted ALOHA for a sensor network]
Five hundred sensors share one radio channel at 10\,kb/s. Each sends a 50-byte report once a minute,
at a random time, to a gateway that acknowledges each one. Can the network use ALOHA?

A packet lasts $400/10^4=40$\,ms, so there are 1500 packet times per minute, and the offered load of
new packets is $S=500/1500=0.333$ packets per packet time. Pure ALOHA cannot carry it at all, since
$S>1/(2e)=0.184$. Slotted ALOHA can: solving $Ge^{-G}=0.333$ on the stable branch gives $G=0.62$, so
each report is transmitted $G/S\approx1.86$ times on average, and each attempt succeeds with probability
$e^{-0.62}=0.54$. The sensors spend 86\% more transmit energy than they would in a collision-free schedule,
and the delay is a geometric number of backoff intervals.

The network saturates at $S=1/e$, that is $1500/e\approx550$ sensors; adding the 551st does not degrade
it gently but risks pushing it across the instability. With pure ALOHA the limit is about 275 sensors,
which is the reason long-range IoT networks such as LoRaWAN, which use unslotted ALOHA, rely on many
parallel channels and spreading factors and on the capture effect (a strong packet often survives a
collision with a weak one) to scale.
\end{worked}

\subsection{Carrier sense: listen before talk}

ALOHA terminals are deaf: they transmit without checking whether the channel is busy. Polite people
do better. Before speaking in a meeting, you listen for a pause. If the propagation delay $\tau$
between terminals is small compared with the packet time $T_p$, a terminal can \emph{listen first} and
defer if it hears another transmission. This is \term{carrier-sense multiple access} (CSMA). Collisions
can now happen only when two terminals start within $\tau$ of each other, before either can hear the
other (Figure~\ref{fig:ch20_csma_window}). The key parameter is the normalised delay
\begin{equation}
a=\frac{\tau}{T_p}.
\end{equation}

\mdcfig{ch20_csma_window}{Why carrier sense cannot prevent every collision. A's signal needs the
propagation delay $\tau$ to reach B; if B senses the channel during that window, it hears silence and
transmits too. Any station that senses after $\tau$ hears A and defers. The smaller $\tau$ is compared
with the packet length, the rarer such collisions become.}

Kleinrock and Tobagi analysed the main variants in 1975.
\begin{itemize}
\item \textbf{Nonpersistent CSMA}: if the channel is busy, wait a random time and sense again.
\item \textbf{1-persistent CSMA}: if busy, keep listening and transmit as soon as the channel goes idle.
This minimises idle time but guarantees a collision whenever two terminals are waiting for the same
transmission to end---like two people who both pounce the instant a speaker pauses.
\item \textbf{$p$-persistent CSMA} (slotted): when the channel goes idle, transmit with probability $p$, or
defer one slot with probability $1-p$, a compromise between the two.
\end{itemize}

\mdcfig{ch20_aloha}{Random-access throughput. (a) Throughput versus offered load (including
retransmissions) for pure and slotted ALOHA and three CSMA variants with normalised propagation delay
$a=0.01$, plus nonpersistent CSMA with $a=0.1$. (b) Maximum throughput as a function of $a$: carrier
sensing helps only when the propagation delay is a small fraction of the packet time. The CSMA/CD curve
is the Metcalfe--Boggs efficiency $1/(1+2(e-1)a)$.}

The throughput of nonpersistent CSMA follows from a renewal argument, sketched in the box below:
\begin{equation}
S_{\text{np}}=\frac{G\,e^{-aG}}{G(1+2a)+e^{-aG}} .
\label{eq:ch20:csma}
\end{equation}
As $a\to0$ this tends to $G/(1+G)$, which approaches 1 at high load: with instantaneous sensing,
collisions vanish. Figure~\ref{fig:ch20_aloha} compares the variants. For $a=0.01$, nonpersistent CSMA
reaches about 0.82 and its slotted version about 0.86, while 1-persistent CSMA peaks near 0.53 because of
its built-in collisions. Panel (b) shows the essential fact about carrier sensing: its benefit depends
entirely on $a$. When the propagation delay approaches the packet time, as on a satellite channel with
$\tau\approx0.25$\,s, carrier sense is useless and ALOHA is as good as anything: by the time you hear
somebody else, their whole packet is already on its way.

\begin{deeper}[The renewal argument for nonpersistent CSMA]
The channel alternates between idle periods and busy periods. With Poisson attempts at rate $G$ per
packet time, an idle period lasts $1/G$ on average. A busy period starts with a packet; any packet that
starts within the next $a$ (before the first is heard) collides with it, and the busy period ends $1+a$
after the last such packet starts. The busy period is successful with probability $e^{-aG}$ (no other
start within $a$) and has mean length $1+2a-(1-e^{-aG})/G$. Dividing the mean useful time per cycle,
$e^{-aG}$, by the mean cycle length, $1/G+1+2a-(1-e^{-aG})/G$, gives~\eqref{eq:ch20:csma}.
\end{deeper}

\subsection{CSMA/CD and Ethernet}

On a cable a station can do more than listen before transmitting: it can listen \emph{while}
transmitting, detect a collision within one round-trip time, and abort---the polite speaker who stops
mid-word on hearing someone else start. This is \term{CSMA with collision detection} (CSMA/CD), the
protocol of the original Ethernet. A collided transmission now wastes only about $2\tau$ instead of a
whole packet.

\mdcpairany[htbp]{ch20_metcalfe_memo}{The first two pages of Metcalfe's memo of 22~May 1973. He proposes to
stop calling the project ``the Alto ALOHA network'' and to call it ``the Ether Network'', and sketches
cable, telephone and radio ``ethers'' joined by boosters.}{ch20_ethernet_tap}{An early 3Com Ethernet unit
(model 3C100), autographed by Bob Metcalfe, who co-founded 3Com in 1979. The two connectors join it to
the shared cable on which every station listened, talked and backed off.}

\begin{history}[Ethernet]
Robert Metcalfe joined Xerox PARC in 1972 after writing a doctoral thesis that drew on the ALOHA system. His
memo of 22~May 1973 proposed connecting PARC's Alto workstations with a passive coaxial cable, the ``ether'',
on which stations would contend ALOHA-style but with carrier sense, collision detection and a backoff that
adapted to the load. With David Boggs he built a 2.94\,Mb/s experimental network, and their 1976 paper in
\emph{Communications of the ACM} gave the efficiency analysis that follows. The DEC--Intel--Xerox
specification of 1980 raised the rate to 10\,Mb/s, and IEEE~802.3 standardised it in 1983. Ethernet's
rivals---Token Ring, Token Bus, FDDI---offered deterministic access, and for a decade their advocates argued
that CSMA/CD would collapse under load. Ethernet won on simplicity and cost, and then, with switched
full-duplex links in the 1990s, abolished its own access problem: a modern Ethernet link has exactly one
transmitter in each direction and no collisions at all. The 802.3 standard still describes CSMA/CD, but no
link faster than 1\,Gb/s implements it.
\end{history}

\chTwentyFigH[0.5]{ch20_backoff}{Binary exponential backoff. After each collision the window from which the
random wait is drawn doubles, up to 1024 slots after the tenth; the dots show one random draw per stage.}

Metcalfe and Boggs's analysis is a model of engineering approximation. After each transmission, the
channel enters a contention interval divided into slots of length $2\tau$, the time needed to detect a
collision. Suppose $k$ stations are ready and each transmits in a contention slot with probability $1/k$
(an ideal adaptive backoff). The probability that exactly one transmits, and so acquires the channel, is
\begin{equation}
A=\left(1-\frac1k\right)^{k-1}\;\xrightarrow{\;k\to\infty\;}\;\frac1e,
\end{equation}
and the number of wasted contention slots before a success is geometric with mean $(1-A)/A\to e-1$. The
fraction of time spent carrying frames is therefore
\begin{equation}
\eta=\frac{T_p}{T_p+2\tau(e-1)}=\frac{1}{1+2(e-1)\,a}\approx\frac{1}{1+3.44\,a}.
\label{eq:ch20:csmacd}
\end{equation}
(Other accounting conventions, for example including the jam signal and the propagation of the final
frame, give slightly different constants such as $1/(1+5a)$; the dependence on $a$ is the point.) In words:
each frame pays a contention tax of about 3.4 propagation delays; long frames on short cables pay
almost nothing.

\begin{worked}[Why Ethernet has a minimum frame size]
Classic 10\,Mb/s Ethernet allowed a network diameter of 2500\,m with up to four repeaters and set the
\emph{slot time}, the maximum round-trip time including repeater and transceiver delays, at 512 bit times,
51.2\,$\mu$s. A station must still be transmitting when news of a collision at the far end returns, or it
will not know its frame collided. The minimum frame is therefore 512 bits, 64 bytes, which is why Ethernet
pads short frames.

For a maximum-size frame of 1518 bytes, $T_p\approx1.21$\,ms and, taking $2\tau$ equal to the slot time,
$a\approx0.021$; \eqref{eq:ch20:csmacd} gives an efficiency of about 93\%. For minimum-size frames,
$T_p=51.2\,\mu$s and $a=0.5$, so the efficiency falls to about 37\%. At 1\,Gb/s the same 512-bit slot
would shrink the allowed diameter to about 20\,m, so half-duplex Gigabit Ethernet extended the slot to
4096 bits with \emph{carrier extension}. In practice nobody used it: switched full-duplex links had made
collision detection obsolete.
\end{worked}

Ethernet's retransmission rule is \term{truncated binary exponential backoff}: after the $i$th collision of a
frame, a station waits a random number of slot times chosen uniformly from $0$ to $2^{\min(i,10)}-1$, and
gives up after 16 attempts (Figure~\ref{fig:ch20_backoff}). Each collision is evidence of congestion, and
doubling the window halves the attempt rate, an approximation to the $q_r=1/n$ rule that requires no
estimate of $n$. It is what people do in a crowded doorway: bump once, wait a moment; bump again, wait longer.

\mdcfig{ch20_wall}{Two people on either side of a wall, both talking to the same person in the doorway:
each thinks the room is quiet, and the listener hears a jumble. That is the hidden terminal (a). The
exposed terminal (b) is the opposite mistake: C overhears B and politely keeps quiet, although its
conversation with D could not possibly disturb A.}

\subsection{Wireless: hidden and exposed terminals}

Collision detection does not work on a radio channel. A transmitting radio's own signal is 100\,dB or more
stronger than anything it might hear from other stations, so it cannot listen while it talks (the
full-duplex problem of Section~\ref{sec:ch20:duplex}). Worse, carrier sense on a radio channel answers the
wrong question. It tells the \emph{transmitter} whether the channel is busy near the transmitter, whereas
collisions happen at the \emph{receiver}. Two configurations expose the gap
(Figures~\ref{fig:ch20_wall} and~\ref{fig:ch20_hidden}).
\begin{itemize}
\item The \term{hidden terminal}: A and C both send to B, but cannot hear each other (they are too far
apart, or a wall stands between them). Each senses an idle channel, both transmit, and their packets
collide at B. Carrier sense does not help at all; the pair behaves like ALOHA.
\item The \term{exposed terminal}: B sends to A, and C, which can hear B, wants to send to D on the other
side. C's transmission would not disturb A, which is out of C's range, yet C hears B and defers needlessly,
wasting a spatial reuse opportunity.
\end{itemize}

Tobagi and Kleinrock identified the hidden-terminal problem in 1975 and proposed a busy-tone solution in
which the receiver signals on a separate channel. The modern solution came from amateur packet radio: Phil
Karn's MACA protocol (1990) had the sender transmit a short \term{request to send} (RTS) and the receiver
answer with a \term{clear to send} (CTS). Every station that hears either message learns that a transmission
is coming, and for how long, and keeps quiet. Stations hidden from the sender still hear the receiver's CTS.
A collision can now only destroy a short RTS, not a long data frame. In the wall picture, the person in the
doorway calls out ``go ahead, A---everyone else hold on a minute'', and C, who cannot hear A, can certainly
hear that.

\begin{figure}[htbp]\centering
\begin{tikzpicture}[font=\sffamily\scriptsize,
  st/.style={circle,draw=navy,fill=navy!10,thick,minimum size=0.55cm,inner sep=0pt,font=\sffamily\small\bfseries}]
% geometry: hidden terminal
\begin{scope}
\fill[accent!7] (0,0) circle (1.75); \draw[accent!50,dashed] (0,0) circle (1.75);
\fill[practc!8,opacity=0.8] (3.0,0) circle (1.75); \draw[practc,dashed] (3.0,0) circle (1.75);
\node[st] (A) at (0,0) {A}; \node[st] (B) at (1.5,0) {B}; \node[st] (C) at (3.0,0) {C};
\draw[->,thick,accent] (A) -- (B); \draw[->,thick,practc] (C) -- (B);
\node[accent] at (-0.2,-1.95) {range of A}; \node[practc] at (3.2,-1.95) {range of C};
\node[align=center,font=\sffamily\small\bfseries,text=navy] at (1.5,2.05) {hidden terminal};
\node[align=center] at (1.5,-2.45) {A and C cannot hear each other;\\their frames collide at B};
\end{scope}
% timeline: RTS/CTS and NAV
\begin{scope}[xshift=7.2cm,yshift=1.1cm]
\node[font=\sffamily\small\bfseries,text=navy] at (3.2,0.95) {RTS/CTS and virtual carrier sense};
\foreach \y/\n in {0/A (sender),-0.75/B (receiver),-1.5/C (hidden from A),-2.25/D (hears A only)}
  {\draw[gray] (0,\y) -- (7.0,\y); \node[anchor=east] at (0,\y) {\n};}
\fill[navy!70] (0.2,0.02) rectangle (0.9,0.36); \node[text=white] at (0.55,0.19) {RTS};
\fill[accent!70] (1.1,-0.73) rectangle (1.75,-0.39); \node[text=white] at (1.42,-0.56) {CTS};
\fill[navy!70] (1.95,0.02) rectangle (5.45,0.36); \node[text=white] at (3.7,0.19) {DATA};
\fill[accent!70] (5.65,-0.73) rectangle (6.3,-0.39); \node[text=white] at (5.98,-0.56) {ACK};
\fill[gray!35] (1.75,-1.48) rectangle (6.3,-1.18); \node at (4.0,-1.33) {NAV set from CTS};
\fill[gray!35] (0.9,-2.23) rectangle (6.3,-1.93); \node at (3.6,-2.08) {NAV set from RTS};
\draw[gray,dashed] (6.3,0.45) -- (6.3,-2.4);
\node[align=center] at (3.2,-2.75) {after the ACK: DIFS, then backoff for everyone};
\end{scope}
\end{tikzpicture}
\caption{Left: the hidden-terminal problem. Carrier sense at A and C reports an idle channel while their
frames collide at B. Right: the 802.11 RTS/CTS exchange. Each control frame carries the remaining duration
of the exchange; stations that hear either one set their \emph{network allocation vector} (NAV) and treat
the medium as busy until it expires, so a station hidden from the sender is silenced by the receiver's CTS.}
\label{fig:ch20_hidden}
\end{figure}

\subsection{CSMA/CA and the 802.11 DCF}

The IEEE~802.11 \term{distributed coordination function} (DCF) combines carrier sense, a random backoff
before \emph{every} transmission, acknowledgements and optional RTS/CTS. It is called \term{CSMA with
collision avoidance} (CSMA/CA) because the random backoff is taken before transmitting, in the hope of
avoiding the collision, rather than after detecting one. The rules are:
\begin{enumerate}
\item A station with a frame senses the medium. It must find the medium idle for a \term{DIFS} (DCF
inter-frame space) before it starts counting.
\item It then counts down a backoff counter, chosen uniformly from $0$ to $CW-1$, by one for each idle
slot. If the medium becomes busy, the counter freezes and resumes after the medium has been idle for
another DIFS. When the counter reaches zero, the station transmits.
\item The receiver answers a correct frame with an ACK after a short inter-frame space (SIFS), shorter than
DIFS so that no other station can seize the medium in between.
\item With no ACK, the sender assumes a collision, doubles $CW$ (up to $CW_{\max}$) and retries; after a
success, $CW$ returns to $CW_{\min}$.
\end{enumerate}
Figure~\ref{fig:ch20_dcf_timeline} plays these rules out for three stations. The frozen counters are the
clever part: a station that loses one contention keeps its place in the queue, like a deli customer who
holds on to their ticket number while someone else is served.

\mdcfig{ch20_dcf_timeline}{The 802.11 DCF in action. After the medium has been idle for a DIFS, three
stations count down random backoff counters (5, 3 and 7). Station 2 reaches zero first and transmits;
the others freeze their counters, wait for the SIFS, ACK and another DIFS, then resume where they left
off. Station 1, with the smallest remaining count, wins the next round.}

Carrier sense is both physical (energy and preamble detection) and \emph{virtual}: every frame carries a
duration field, and stations that hear it set a \term{network allocation vector} (NAV) and defer for that
long without sensing. For the OFDM physical layers (802.11a/g/n/ac/ax), the slot is 9\,$\mu$s, SIFS
16\,$\mu$s, DIFS 34\,$\mu$s, $CW_{\min}=16$ and $CW_{\max}=1024$ (the standard writes these as 15 and
1023, the largest counter values). The 802.11e enhanced distributed channel access (EDCA) gives four access
categories---voice, video, best effort, background---different inter-frame spaces and windows, so that, for
example, voice frames use windows of 4--8 slots and win most contentions against bulk data with 16--1024.

\subsection{Bianchi's model of the DCF}

How efficient is the DCF with $n$ stations that always have frames to send? In 2000 Giuseppe Bianchi gave
an answer so accurate and so simple that it became the basis of nearly all subsequent Wi-Fi analysis. Its
key assumption is that every transmission attempt, whatever the station's backoff history, collides with the
same constant probability $p$, independently of everything else. Under this \emph{decoupling} approximation,
each station's backoff process becomes a Markov chain over its backoff stage $i$ (the number of
collisions suffered by the current frame) and its counter value. The intuition is a crowd in which
each person reacts only to the average noise of the room, not to any particular neighbour: the
crowd's behaviour then follows from one number, the probability that a given attempt collides.

\mdcpairany[htbp]{ch20_wifi_ap}{An enterprise Wi-Fi access point (Cisco Aironet 1131AG) of the
802.11a/g generation analysed in Figure~\ref{fig:ch20_bianchi}. Every station associated with it
contends for the same channel using the DCF.}{ch20_airtime}{Where the airtime goes for one 1500-byte
frame from a single 802.11a/g station: inter-frame spaces, the mean backoff, the PHY preamble and
headers and the ACK are fixed in time, so they eat a larger share as the data rate rises.}

Solving the chain for its stationary distribution gives the probability $\tau$ that a station transmits
in a randomly chosen slot:
\begin{equation}
\tau=\frac{2(1-2p)}{(1-2p)(W+1)+pW\left(1-(2p)^m\right)},
\label{eq:ch20:bianchitau}
\end{equation}
where $W=CW_{\min}$ and $m$ is the number of doublings ($CW_{\max}=2^mW$). As a check, with no collisions
($p=0$) this is $2/(W+1)$, one transmission per mean backoff of $(W-1)/2$ slots. A transmission collides if
any of the other $n-1$ stations transmits in the same slot:
\begin{equation}
p=1-(1-\tau)^{n-1}.
\label{eq:ch20:bianchip}
\end{equation}
Equations \eqref{eq:ch20:bianchitau} and \eqref{eq:ch20:bianchip} are two equations in two unknowns, with a
unique solution that is easy to find numerically: the attempt rate sets the collision probability, and the
collision probability sets the attempt rate, and only one pair is self-consistent.

\mdcfig{ch20_bianchi}{Bianchi's model of the 802.11 DCF for an 802.11a/g network at 54\,Mb/s with
1500-byte frames (control frames at 24\,Mb/s). (a) Saturation throughput versus the number of stations:
basic access (DATA--ACK) and RTS/CTS, for $CW_{\min}=16$ and 32; circles are a slot-level simulation of
binary exponential backoff. (b) The conditional collision probability $p$ and the per-slot attempt
probability $\tau$ (scaled by 5) from the fixed point.}

\begin{deeper}[From the fixed point to throughput]
The saturation throughput follows from counting what happens in a generic slot. A slot contains a
transmission with probability $P_{\text{tr}}=1-(1-\tau)^n$; given a transmission, it is a success with
probability $P_s=n\tau(1-\tau)^{n-1}/P_{\text{tr}}$. If $\sigma$ is the empty-slot duration, $T_s$ and $T_c$
the durations of a successful and a collided exchange, and $T_{\text{pl}}$ the airtime of the payload,
\begin{equation}
S=\frac{P_s P_{\text{tr}}\,T_{\text{pl}}}{(1-P_{\text{tr}})\,\sigma+P_{\text{tr}}P_s\,T_s+P_{\text{tr}}(1-P_s)\,T_c}.
\label{eq:ch20:bianchiS}
\end{equation}
The numerator is the expected useful airtime per slot; the denominator is the expected length of a slot,
which is $\sigma$ when nobody transmits and stretches to a whole exchange when somebody does.
\end{deeper}

Figure~\ref{fig:ch20_bianchi} evaluates the model for an 802.11a/g network at 54\,Mb/s and compares it with a
slot-level simulation of the actual backoff rules; the agreement is within a few percent. Several practical
lessons emerge.
\begin{itemize}
\item Even a single station gets only about 30\,Mb/s of a 54\,Mb/s link: preambles, inter-frame spaces,
ACKs and the mean backoff consume the rest (Figure~\ref{fig:ch20_airtime}). This overhead is fixed in time,
so it hurts more as the PHY rate rises, and it is the main reason 802.11n and later standards aggregate many
frames into one transmission (A-MPDU) and acknowledge them with one block ACK.
\item With basic access, throughput falls steadily with $n$, by nearly 30\% from 2 to 50 stations, as the
collision probability climbs above 50\%. With RTS/CTS, collisions cost only an RTS, so throughput is almost
independent of $n$, but every success pays for the extra handshake; the curves cross near 20 stations for
these parameters. This is why RTS/CTS is usually enabled only for large frames, through an \emph{RTS threshold}.
\item A larger $CW_{\min}$ wastes idle slots with few stations but reduces collisions with many. No fixed
window is optimal for all $n$; Bianchi showed that the optimal attempt probability is roughly
$\tau\approx\frac1n\sqrt{2\sigma/T_c}$, so the ideal window should grow with the number of stations. Wi-Fi~6
and~7 sidestep the problem in dense networks by letting the access point schedule uplink OFDMA transmissions
with trigger frames, turning contention into reservation.
\end{itemize}

\begin{tryit}[Crowd a Wi-Fi channel]
In Lab 26's \emph{Wi-Fi's DCF: Bianchi's model}, drag \emph{Saturated stations} from 1 to 60 with
\emph{Access} set to Basic, then switch to RTS/CTS and find the crossover. Change \emph{Minimum contention
window} from 16 to 64 and watch the curve tilt: worse with few stations, better with many. For a live
version, open Lab 11's \emph{Wi-Fi contention, live}, raise \emph{Saturated stations} and toggle
\emph{RTS/CTS handshake} to see collisions shrink from data-sized to RTS-sized.
\end{tryit}

\subsection{Reservation, polling and hybrids}

Between fixed assignment and pure contention lies a family of hybrid protocols that use a little contention
to make reservations and then transmit without collisions---the restaurant where you queue once for a table
and then eat in peace.

\chTwentyFigH{ch20_reservation}{Request--grant access, the pattern of DOCSIS upstreams, satellite return links
and, in spirit, LTE and NR scheduling requests. Short request minislots are contended ALOHA-style and
occasionally collide; the long data slots that follow are granted by the controller and never do.}

\begin{itemize}
\item \textbf{Polling.} A central station invites each terminal in turn. Bluetooth's central device polls
its peripherals; 802.11's point coordination function (PCF) and the 802.11e HCCA did the same, though they were
little used. Polling is collision-free and gives bounded delay, but wastes time polling idle terminals.
\item \textbf{Request--grant.} Terminals send short requests in contention mini-slots and the controller
grants dedicated transmission slots. Cable modems (DOCSIS) work this way on the upstream, as do satellite
return links with demand-assigned multiple access (Chapter~\ref{ch:satellite}), and the scheduling requests
of LTE and NR follow the same pattern with a dedicated rather than contended request channel.
\item \textbf{Reservation ALOHA and PRMA.} A terminal that wins a slot in contention keeps it in subsequent
frames for as long as it has data. Goodman's packet reservation multiple access (PRMA, 1989) applied this to
voice with silence suppression: a talkspurt reserves a slot that is released during silence.
\item \textbf{Trigger-based uplink OFDMA.} In 802.11ax the access point sends a trigger frame assigning
resource units to stations, some of which can be left for random access (UL OFDMA-based random access),
a slotted ALOHA within an OFDMA grid.
\end{itemize}
Cellular systems use random access only to get a terminal's foot in the door. Everything after that is
scheduled, which is the subject of the next section and of Section~\ref{sec:ch20:scheduling}.

%======================================================================
\section{Random access in cellular systems}
\label{sec:ch20:rach}
%======================================================================

Switch on a phone, or let it sit idle in your pocket for a few seconds, and it has no resources at all.
Before it can send a single byte it must knock on the door. A phone that has just been switched on, or that
has been idle and now has data, or that has lost uplink synchronisation, must announce itself on the
\term{random-access channel} (RACH). The cellular RACH is a slotted ALOHA protocol with three refinements:
transmissions are designed to be detectable even when they collide, the base station uses the first message
to measure the terminal's timing, and contention is resolved explicitly a few messages later.

\mdcfig{ch20_zc}{Detecting preambles. Two terminals send different cyclic shifts of the same 839-long
Zadoff--Chu root with different round-trip delays, buried in noise. One circular correlation (an FFT,
a multiplication and an inverse FFT) produces a peak in each preamble's zone of $N_{CS}$ lags: the zone
identifies the preamble, and the offset within the zone measures the delay. The weaker terminal is
detected as easily as the stronger.}

\subsection{Preambles}

In LTE and 5G NR, the terminal's first transmission is a \term{preamble}, one of (typically) 64 sequences
available in the cell, transmitted on the physical random-access channel (PRACH) in periodically recurring
\term{RACH occasions}. The preambles are cyclic shifts of Zadoff--Chu sequences (Chapter~\ref{ch:sync}) of
length 839 (for the long formats, at 1.25 or 5\,kHz subcarrier spacing) or 139 (for NR's short formats).
Zadoff--Chu sequences have ideal periodic autocorrelation: every cyclic shift of a root sequence is orthogonal
to every other. The base station correlates the received signal against each root sequence with one FFT, and
the position of a correlation peak tells it both \emph{which} preamble was sent and \emph{when} it arrived,
giving the round-trip delay from which it computes the terminal's timing advance
(Figure~\ref{fig:ch20_zc}). The cyclic-shift spacing must exceed the largest round-trip delay plus the delay
spread, otherwise a late-arriving preamble looks like a different one; large cells therefore get fewer shifts
per root and use more roots.

Several terminals that choose \emph{different} preambles in the same occasion do not collide at all: the base
station detects them all, as in CDMA. Only terminals that choose the \emph{same} preamble collide, and the
collision is not even visible at this stage, because the base station sees a single preamble (perhaps with two
delay peaks). The RACH is therefore a slotted ALOHA with $M$ parallel ``channels'' per occasion, where $M$ is
the number of preambles available for contention, often 54 of the 64, with the rest reserved for
contention-free access. If $k$ terminals contend, each picking a preamble uniformly at random, the expected
number that pick a preamble nobody else picked is
\begin{equation}
\bar N_{\text{succ}}(k)=k\left(1-\frac1M\right)^{k-1}\approx k\,e^{-k/M},
\label{eq:ch20:rach}
\end{equation}
the slotted-ALOHA law again, with load $G=k/M$. It peaks at $k=M$ with about $M/e$ successes per occasion
(Figure~\ref{fig:ch20_rach}a). It is the birthday problem in disguise: with 54 ``birthdays'' to choose from,
a room of only nine people already has better-than-even odds that two of them share one.

\begin{figure}[htbp]\centering
\begin{tikzpicture}[x=1.15cm,font=\sffamily\scriptsize,>={Stealth[length=2mm]}]
\begin{scope}
\node[font=\sffamily\small\bfseries,text=navy] at (2,0.75) {four-step (LTE, NR)};
\node[draw=navy,fill=navy!8,thick,minimum width=1.1cm] at (0,0.2) {UE};
\node[draw=navy,fill=navy!8,thick,minimum width=1.1cm] at (4,0.2) {gNB};
\draw[gray] (0,-0.1) -- (0,-4.6); \draw[gray] (4,-0.1) -- (4,-4.6);
\draw[->,thick,accent] (0,-0.5) -- node[above,sloped]{Msg1: preamble on PRACH} (4,-1.0);
\draw[->,thick,navy] (4,-1.4) -- node[above,sloped]{Msg2: RAR (TA, grant, temp.\ ID)} (0,-1.9);
\draw[->,thick,accent] (0,-2.3) -- node[above,sloped]{Msg3: identity, RRC request} (4,-2.8);
\draw[->,thick,navy] (4,-3.2) -- node[above,sloped]{Msg4: contention resolution} (0,-3.7);
\node[align=left,anchor=west,text=gray] at (4.15,-2.75) {collisions\\surface here};
\node[align=center,text=gray] at (2,-4.25) {two round trips before data};
\end{scope}
\begin{scope}[xshift=7.6cm]
\node[font=\sffamily\small\bfseries,text=navy] at (2,0.75) {two-step (NR Release 16)};
\node[draw=navy,fill=navy!8,thick,minimum width=1.1cm] at (0,0.2) {UE};
\node[draw=navy,fill=navy!8,thick,minimum width=1.1cm] at (4,0.2) {gNB};
\draw[gray] (0,-0.1) -- (0,-4.6); \draw[gray] (4,-0.1) -- (4,-4.6);
\draw[->,thick,accent] (0,-0.5) -- node[above,sloped]{MsgA: preamble + PUSCH payload} (4,-1.0);
\draw[->,thick,navy] (4,-1.4) -- node[above,sloped]{MsgB: response + resolution} (0,-1.9);
\node[align=center,text=gray] at (2,-2.9) {one round trip; falls back\\to four-step if only the\\preamble is decoded};
\end{scope}
\end{tikzpicture}
\caption{Contention-based random access. In the four-step procedure, terminals that pick the same preamble
discover the collision only after Msg3; the two-step procedure sends the preamble and the payload together and
saves a round trip at the cost of reserving uplink resources for payloads that may collide.}
\label{fig:ch20_rachsteps}
\end{figure}

\subsection{The four-step procedure}

Contention-based random access in LTE and NR proceeds in four messages (Figure~\ref{fig:ch20_rachsteps}).
\begin{enumerate}
\item \textbf{Msg1, preamble.} The terminal picks a preamble at random and transmits it in the next RACH
occasion, at a power computed from its downlink path-loss estimate.
\item \textbf{Msg2, random-access response} (RAR). The base station answers, on the downlink shared channel,
for each detected preamble: a timing-advance command, an uplink grant for the next message, and a temporary
identifier. A terminal that receives no response within its window retransmits the preamble with a power
increased by a configured \emph{ramping step}, after a random backoff that the base station can lengthen
through a backoff indicator when the channel is overloaded.
\item \textbf{Msg3, scheduled transmission.} The terminal, now time-aligned, sends its identity and the reason
for access (for example an RRC connection or resume request) on the granted resource. If two terminals chose
the same preamble, both received the same RAR and both transmit Msg3 on the same resource. This is where the
collision actually happens.
\item \textbf{Msg4, contention resolution.} The base station echoes the identity it decoded from Msg3. The
terminal whose identity matches has won; the others restart.
\end{enumerate}
\term{Contention-free random access} skips the lottery: for handover, beam-failure recovery or when the network
wants a terminal to synchronise, the base station assigns it a dedicated preamble in advance, and steps 3 and 4
are unnecessary.

The two-step procedure, added in NR Release~16, merges Msg1 with Msg3 (``MsgA'': a preamble followed by a
payload on a pre-configured uplink resource) and Msg2 with Msg4 (``MsgB''). It halves the access latency,
which matters for small cells, for unlicensed spectrum (where each message must win a listen-before-talk
contention) and for low-latency services. Its cost is that payload resources are reserved in advance and wasted
when unused or collided. If the base station detects the preamble but cannot decode the payload, it can fall
back to the four-step procedure by answering with an ordinary RAR.

\mdcfig{ch20_rach}{RACH contention. (a) Expected successful and collided devices in one RACH occasion versus
the number of contenders, for 54 and 64 preambles; successes peak at about $M/e$ when $k\approx M$.
(b) Simulated access success probability (at most 10 preamble attempts, uniform backoff up to 20\,ms) with
54 preambles and an occasion every 5\,ms, for devices activating uniformly over 60\,s or in a Beta(3,4) burst
over 10\,s. Our simplified simulation ignores limits on Msg2 and Msg3 resources, which in a real network cause
overload even earlier.}

\subsection{Massive access}

Machine-type communication changes the RACH problem. A cell may serve tens of thousands of meters, sensors and
trackers, each sending tiny amounts of data, and their activity is sometimes correlated: a power outage and
restoration, or an alarm, can make thousands of devices attempt access within seconds. 3GPP's study of this
scenario (TR~37.868) defined a ``bursty'' traffic model in which up to 30\,000 devices in a cell activate over
10\,s with a Beta(3,4) distribution in time, as opposed to a uniform spread over 60\,s.

Figure~\ref{fig:ch20_rach}b shows what happens. With 54 preambles and a RACH occasion every 5\,ms, the cell
can admit at most about $20$ devices per occasion, or roughly 4000 per second. A uniform 60-second spread of
30\,000 devices (500 per second) is easily handled. The same devices in a 10-second burst peak at about
6000 arrivals per second; the RACH tips into the collapsed region of the ALOHA curve, retries compound the
overload, and in our simplified simulation only about a third of the devices get through. The defences are
the ones the theory suggests:
\begin{itemize}
\item \textbf{Spreading the load in time}: \emph{access class barring} (ACB), in which the base station broadcasts
a barring probability and a barring time so that devices admit themselves with a controlled probability (an
explicit version of $q_r=1/\hat n$); and \emph{extended access barring} (EAB), which bars delay-tolerant devices
entirely during overload.
\item \textbf{More resources}: more RACH occasions or preambles, or separate RACH resources for machines so that
they cannot block phones.
\item \textbf{Fewer messages}: early data transmission in Msg3 (LTE-M and NB-IoT, Release~15), two-step access,
and configured-grant or grant-free uplink, which avoid random access altogether for periodic traffic.
\end{itemize}
NB-IoT adds a twist: its single-tone, frequency-hopping NPRACH and its \emph{coverage-enhancement levels} let a
device in a basement repeat its preamble many times, so devices with poor coverage occupy the channel for longer.

\chTwentyPhotoH[0.5]{ch20_crowd}{Wembley Stadium, London, as a crowd of tens of thousands streams along the
approach. When such a crowd arrives or leaves, the phones in it wake up together, and the random-access
channel, not the data capacity, is often the first thing to saturate.}

\begin{tryit}[Flood the RACH, then bar it]
Open Lab 26's \emph{RACH collisions and barring}. First drag \emph{Devices contending} past
\emph{Preambles} and watch successes peak near $M/e$. Then set \emph{Devices} to 30\,000 with
\emph{Activation} at 10\,s burst: most devices fail. Now lower \emph{Access-class barring: pass probability}
step by step and find the value that lets the most devices through---slowing everyone down gets more of
them in.
\end{tryit}

\begin{pitfall}[Random access is where networks fail first]
Engineers dimension networks for traffic, but outages often begin with signalling. When a large venue empties,
when a network recovers after a fault, or when a firmware update makes a whole fleet of devices reconnect at the
same moment, the RACH and the control channels behind it saturate long before the data channels do. Because
ALOHA-like systems have a stable congested state (Figure~\ref{fig:ch20_aloha_stability}a), the overload can
persist after its cause has gone. Device makers are therefore told to randomise their reconnection times, and
operators keep overload controls (barring, backoff indicators, rejection with wait timers) configured and tested.
A fleet of devices that all report ``at midnight'' is a self-inflicted denial-of-service attack.
\end{pitfall}

%======================================================================
\section{The cellular concept}
\label{sec:ch20:cellular}
%======================================================================

In 1976 New York had twelve mobile channels and a waiting list of thousands. Today a single city block
may host more simultaneous mobile connections than the whole of North America had subscribers then.
The spectrum grew, and the radios got better, but neither explains most of the difference. The idea
that does was written down in a Bell Labs memo three years after the first mobile call.

\chTwentyFigH[0.45]{ch20_capacity_cellsize}{Equation~\eqref{eq:ch20:capacity} for a 100\,km$^2$ city with the
395 AMPS voice channels of one operator and $N=7$: every halving of the cell radius quadruples the number
of simultaneous calls, from about two thousand with 1\,km cells to more than two hundred thousand
with 100\,m cells.}

\begin{history}[From a 1947 memo to a 1983 launch]
On 11~December 1947 Douglas H.~Ring of Bell Laboratories wrote a technical memorandum, ``Mobile
Telephony---Wide Area Coverage'', building on ideas of his colleague W.~Rae Young. It proposed
replacing the single powerful transmitter of a metropolitan mobile service by a lattice of
low-power transmitters, each serving a small area, with the same frequencies reused in areas far
enough apart, and it sketched hexagonal cells. The concept was sound, but it needed three things
that 1947 did not have: spectrum, switching computers that could track a moving call from cell to
cell, and cheap radios. The spectrum took longest. The FCC opened Docket 18262 in 1968 to consider
reallocating the upper UHF television channels to land mobile use; it proposed 75\,MHz near
850\,MHz in 1970, allocated 40\,MHz in 1974, and only in 1981 adopted final rules that split the
cellular band in each market into two 20\,MHz licences, one for the local wireline telephone company
and one for a competitor (enlarged to 25\,MHz each in 1986). Meanwhile AT\&T's Bell Labs designed the
Advanced Mobile Phone System (AMPS), described in a celebrated January 1979 issue of the \emph{Bell
System Technical Journal}, and Illinois Bell ran a developmental AMPS system in Chicago from 1978 with
about two thousand customers. Motorola, competing for the same future, had made the point in public
on 3~April 1973, when Martin Cooper called a rival at Bell Labs from a prototype handheld DynaTAC on a
Manhattan sidewalk. Commercial AMPS service opened in Chicago in October 1983. NTT had launched a
cellular service in Tokyo four years earlier, and the Nordic NMT system followed in 1981.
\end{history}

\subsection{Why cells multiply capacity}

Suppose an operator holds $S$ duplex channels. A single high tower serving a city can carry $S$
simultaneous calls, however large the city. Divide the city instead into cells, partition the channels
into $N$ disjoint sets of $S/N$ channels, and assign the sets so that cells using the same set, called
\term{co-channel cells}, are far apart. A group of $N$ cells that together use all $S$ channels is a
\term{cluster}, and $N$ is the \term{cluster size} or \term{reuse factor}. If the service area is covered
by $M$ clusters, the system carries
\begin{equation}
C=M\cdot N\cdot\frac SN=M\,S
\label{eq:ch20:capacity}
\end{equation}
\chTwentyFigH[0.36]{ch20_map_colours}{Frequency planning is map colouring. Here an irregular network of cells
is coloured so that no cell shares a channel set with a neighbour or a neighbour's neighbour; real cell
layouts, unlike hexagons, need extra sets and careful hand-tuning.}
\noindent simultaneous calls. Capacity is proportional to the number of clusters, so halving the cell radius
quadruples it (Figure~\ref{fig:ch20_capacity_cellsize}). This is the whole idea: \emph{capacity is created
by reuse, and reuse is limited by interference}. The two design questions are how small $N$ can be made,
which is a question of co-channel interference, and how small the cells can be made, which is a question
of cost.

\begin{analogy}[Paint colours on a map]
Frequency reuse is like colouring a map so that no two neighbouring countries share a colour. You do not
need a different colour for every country---a handful suffices, reused across the map. Channel sets are
the colours, and the radio rule is stricter than the cartographer's: two cells with the same set must be
separated by several others, not just by a border, because radio does not stop at boundaries. Fewer
colours (a smaller cluster) mean more channels per cell but nearer same-coloured neighbours, which is
exactly the trade-off the rest of this section quantifies (Figure~\ref{fig:ch20_map_colours}).
\end{analogy}

\chTwentyFigH[0.98]{ch20_hex_clusters}{Hexagonal reuse patterns for cluster sizes 3, 4 and 7. Cells with the
same number use the same channel set; the dark cells are the co-channel cells of channel set 1. The arrows
show the shift rule: move $i$ cells in one direction, turn $60^\circ$ and move $j$ cells to reach the
nearest co-channel cell.}

\subsection{Hexagonal geometry and the cluster size}

Real cells are irregular regions shaped by terrain and buildings, but planning starts from an idealised
layout of equal cells. Of the three regular polygons that tile the plane---triangles, squares and
hexagons---the hexagon is closest to the circle a base station would cover in flat terrain, so it
covers a given area with the fewest cells (bees discovered the same thing for honeycomb). A hexagonal
cell of circumradius $R$ (centre to vertex) has area $\frac{3\sqrt3}{2}R^2$, and the centres of adjacent
cells are $\sqrt3R$ apart.

Which cluster sizes tile the plane so that every cell has its co-channel cells arranged symmetrically
around it? Locate the nearest co-channel cell by walking from the centre of a cell across $i$ cells
in a straight line, turning $60^\circ$, and walking $j$ more (Figure~\ref{fig:ch20_hex_clusters}). The
two legs have lengths $i\sqrt3R$ and $j\sqrt3R$ and meet at $120^\circ$, so by the law of cosines the
\term{co-channel reuse distance} $D$ satisfies
\begin{equation}
D^2=3R^2\left(i^2+j^2+2ij\cos60^\circ\right)=3R^2\left(i^2+ij+j^2\right).
\end{equation}
The co-channel cells themselves form a hexagonal lattice with spacing $D$, and each one ``owns'' a
hexagon of area proportional to $D^2$, made up of the $N$ cells of its cluster. The ratio of areas gives
the cluster size:
\begin{equation}
N=\frac{D^2}{3R^2}=i^2+ij+j^2,\qquad
Q\equiv\frac DR=\sqrt{3N}.
\label{eq:ch20:cluster}
\end{equation}
The allowed cluster sizes are therefore $N=1,3,4,7,9,12,13,16,19,21,\dots$; no layout gives a symmetric
reuse pattern with $N=5$ or $N=6$. The \term{co-channel reuse ratio} $Q=D/R$ is the single number that
controls interference: it says how many cell radii separate you from the nearest stranger on your channel.

\subsection{Co-channel interference}

Consider a terminal at the edge of its cell, at distance $R$ from its own base station, in an
interference-limited system (thermal noise negligible). With the log-distance path-loss model of
Chapter~\ref{ch:channels}, received power falls as $d^{-n}$, where $n$ is the path-loss exponent. If
there are $i_0$ co-channel interferers at distances $D_k$, all transmitting at the same power, the
signal-to-interference ratio is
\begin{equation}
\frac SI=\frac{R^{-n}}{\sum_{k=1}^{i_0}D_k^{-n}} .
\label{eq:ch20:sirgeneral}
\end{equation}
A hexagonal layout has six co-channel cells in the first tier. Approximating all their distances by $D$
and ignoring the second and later tiers (six cells at $\sqrt3D$ and six at $2D$, which add about 17\%, or 0.7\,dB,
to the interference for $n=4$),
\begin{equation}
\frac SI\approx\frac{(D/R)^n}{6}=\frac{\left(\sqrt{3N}\right)^{n}}{6}.
\label{eq:ch20:sir}
\end{equation}
In words: SIR grows as the reuse ratio to the power of the path-loss exponent. For $n=4$,
$S/I\approx\frac32N^2$: doubling the cluster size gains 6\,dB of SIR and halves the channels per cell.
A more cautious estimate places the terminal at the worst corner of its cell, where two interferers are
at about $D-R$, two at $D$ and two at $D+R$:
\begin{equation}
\left.\frac SI\right|_{\text{worst}}\approx\frac{1}{2(Q-1)^{-n}+2Q^{-n}+2(Q+1)^{-n}} .
\label{eq:ch20:sirworst}
\end{equation}
Figure~\ref{fig:ch20_sir_reuse}a plots both. The worst case is 1--3\,dB below the simple estimate for the
cluster sizes of practical interest.

\mdcfig{ch20_sir_reuse}{(a) First-tier co-channel SIR at the cell edge versus cluster size, from the simple
formula~\eqref{eq:ch20:sir} for three path-loss exponents and from the worst-case
formula~\eqref{eq:ch20:sirworst}; dots mark the valid cluster sizes. (b) Simulated distribution of downlink
SIR over locations in a cell, with two tiers of interferers, $n=4$ and independent 8\,dB log-normal shadowing
on every link, for omnidirectional cells and for three 120$^\circ$ sectors with a 65$^\circ$ beamwidth
antenna.}

\mdcfig{ch20_sir_map}{Downlink SIR across a hexagonal network with path-loss exponent 3.5, no shadowing,
each location served by the nearest base station (white dots). (a) Reuse 1: every cell border is a valley of
SIR near 0\,dB. (b) Reuse 3 (channel sets A, B, C): the valleys fill in, at the cost of using one third of the
spectrum in each cell.}

\begin{worked}[Choosing the cluster size for AMPS]
AMPS used analog FM with 30\,kHz channels. Subjective tests at Bell Labs found that FM voice was judged
good by most listeners when the carrier-to-interference ratio was at least about 18\,dB. With $n=4$, what
cluster size is needed?

From~\eqref{eq:ch20:sir}, $(3N)^2/6\ge10^{1.8}=63.1$ gives $3N\ge19.5$, so $N\ge6.5$, and the smallest
valid cluster is $N=7$ ($i=2$, $j=1$), with $Q=4.58$ and $S/I\approx21^2/6=73.5$, or 18.7\,dB. The
worst-case estimate~\eqref{eq:ch20:sirworst} gives only 17.3\,dB, short of the target. There are two
remedies. The next valid cluster, $N=9$, gives 19.8\,dB worst-case but cuts the channels per cell by 22\%.
Alternatively, keep $N=7$ and split each cell into three $120^\circ$ sectors, each with its own third of the
cell's channels. A sector antenna sees only the two first-tier co-channel sectors that face it, so the simple
estimate becomes $Q^4/2=220$, or 23.4\,dB, and even the worst case exceeds 24\,dB. AMPS networks used $N=7$
with three sectors, giving 21 channel sets. With the 416 channel pairs of one AMPS operator (after the 1986
expansion), of which 21 carried control signalling, each cell had 56 or 57 voice channels.

The simulation of Figure~\ref{fig:ch20_sir_reuse}b adds the reality of 8\,dB shadowing. The 18\,dB target is
then met at only about 60\% of locations in an omnidirectional $N=7$ layout and about 80\% with sectors;
shadowing, not geometry, sets the margins. This is one reason the digital systems that followed AMPS, which
could tolerate a much lower SIR thanks to channel coding, achieved large capacity gains.
\end{worked}

The idealised formula assumes every co-channel cell is transmitting. On the uplink the interferers are
terminals scattered around their own cells, and averaging over their positions gives a similar result. On the
downlink of a system with partially loaded cells, interference scales with the load: a cell at 50\% load
radiates roughly half the interference power, an effect that later systems with power control and
discontinuous transmission exploited heavily.

Figure~\ref{fig:ch20_sir_map} shows the same physics as a map. With reuse 1, every cell edge is a region of
SIR near 0\,dB, and the median over the cell is only about 5\,dB for $n=3.5$; with reuse 3, the median rises to
about 14.5\,dB and the worst 5\% of locations still have more than 8\,dB. A narrowband analog or GSM-like system
must use the right-hand pattern. A modern OFDMA system with powerful codes and link adaptation runs the
left-hand pattern and accepts low rates at the cell edges, because using all the spectrum everywhere more than
compensates. Section~\ref{sec:ch20:interference} shows how the cell edges are then rescued.

\mdcpairany[htbp]{ch20_cell_site}{A lattice cell tower on a university campus. The groups of rectangular
panels at the top are sector antennas pointing in different directions, each one a separate ``cell'' covering
a slice of the compass.}{ch20_sector_pattern}{Three sector antennas with the 3GPP-style pattern used in this
chapter's simulations: a 65$^\circ$ half-power beamwidth and a 20\,dB floor. Each sector is deaf, to
within the floor, to interferers behind it.}

\begin{tryit}[Choose a cluster size]
In Lab 26's \emph{Clusters, reuse and SIR}, start at \emph{Cluster size} $N=7$ with \emph{Antennas} set to
Omni and \emph{Shadowing $\sigma$} at 8\,dB, and read off the fraction of locations below the 18\,dB
\emph{SIR target}. Switch \emph{Antennas} to 3 sectors and watch the CDF jump right; then try $N=4$ and
$N=3$ with sectors. Set $\sigma$ to 0 to see the geometry alone. Lab 11's \emph{Frequency reuse and SIR}
adds the other half of the trade-off: drag \emph{Channels in the system} and watch the channels per cell
fall as $N$ grows.
\end{tryit}

\subsection{Sectorization}

Replacing a base station's omnidirectional antenna by three antennas, each covering a $120^\circ$ sector, or six
covering $60^\circ$, reduces the number of first-tier interferers seen by each sector from six to about two
($120^\circ$) or one ($60^\circ$). It is the cheapest SIR improvement available, about 4.8\,dB for three sectors in
the simple model, and almost every macro cell site in the world today is three-sectored
(Figure~\ref{fig:ch20_cell_site}). The SIR gain can be spent in two ways: to improve quality at a fixed cluster
size, as in the AMPS example, or to reduce the cluster size and so increase the number of channels per cell. With
three sectors, $N=4$ gives $(\sqrt{12})^4/2=72$, or 18.6\,dB, in the simple model, which meets the AMPS target
with 75\% more channels per cell than $N=7$.

\chTwentyFigH[0.38]{ch20_cell_splitting}{Cell splitting where the traffic is: the busiest macro cells (blue) are
subdivided into cells of half the radius (orange), each with its own lower-power base station. Each split
cell reuses the full set of channels in a quarter of the area.}

Sectorization has costs. Each sector is a separate trunk with a third of the channels, which reduces trunking
efficiency (Section~\ref{sec:ch20:traffic}); a terminal moving around the site must hand over between sectors;
and real antennas leak power behind and beside them, with front-to-back ratios of perhaps 20--30\,dB, so the
interference reduction is smaller than the idealised count suggests. Sectors are also cells in their own right:
in LTE and NR, every sector is a separate cell with its own identity.

\subsection{Cell splitting, small cells and heterogeneous networks}

The second lever is the cell size. \term{Cell splitting} subdivides a congested cell into smaller cells, each with
its own base station at reduced height and power (Figure~\ref{fig:ch20_cell_splitting}). If the radius is halved,
the number of cells in an area quadruples, and so does the capacity by~\eqref{eq:ch20:capacity}; to keep the same
received power at the new, closer cell edge, the transmit power can drop by $10n\log_{10}2$\,dB, 12\,dB for $n=4$.
Splitting is done where the traffic is, so real networks have cells of very different sizes: rural macro cells
tens of kilometres across, urban macro cells of a few hundred metres, and small cells within them.

The modern network is a \term{heterogeneous network} (HetNet): a layer of macro cells for coverage and mobility,
overlaid with \term{small cells}---microcells on lamp posts and building walls, picocells indoors, femtocells in
homes---that add capacity where users concentrate. Each small cell reuses the full spectrum in a tiny area, so the
capacity added per small cell is roughly that of a macro cell, at a fraction of the cost. Small cells raise their
own problems: their coverage areas are small because their transmit power is low, so they attract few users unless
the network biases association towards them (Section~\ref{sec:ch20:interference}); the macro layer interferes with
them; and a femtocell with closed access becomes a strong interferer for a passing macro user who is not allowed
to connect to it. Fast-moving users are better kept on the macro layer, an ``umbrella'' over the small cells, to
avoid a storm of handovers.

\mdcphoto[0.62]{ch20_small_cell}{A small cell on a lamp post in Minneapolis, photographed during the week of
Super Bowl LII, when the city's crowds peaked: four compact panel antennas with their radios, mounted a few
metres above the street. Each one adds a full cell's capacity to a few city blocks.}

\begin{keyidea}[Spatial reuse is the dominant source of capacity growth]
Over the history of mobile radio, the capacity per unit area has grown by a factor of perhaps a million. Estimates
attributed to Martin Cooper and widely repeated divide the credit roughly as follows: more spectrum and better
modulation and coding have each contributed factors of tens, while the reuse of spectrum through ever smaller cells
has contributed a factor of a thousand or more. The figures are approximate, but the moral is robust: when an
engineer needs a hundredfold capacity increase, the first place to look is not the modem but the cell size.
\end{keyidea}

%======================================================================
\section{Traffic engineering: Erlang's formulas}
\label{sec:ch20:traffic}
%======================================================================

A cell with $C$ channels does not serve $C$ subscribers. It serves many more, because subscribers use their
phones only occasionally, and the question becomes statistical: how many channels does a cell need so that a call
attempt finds them all busy only rarely? The theory that answers it was created for telephone exchanges by a
Danish mathematician more than a century ago, and it is still used to dimension voice networks, call centres and
signalling resources.

\mdcpairany[htbp]{ch20_ktas}{The Central \O{}bro exchange building of KTAS, the Copenhagen Telephone Company
that employed Erlang, in Copenhagen. Questions about how much switching equipment buildings like this
needed gave birth to teletraffic theory.}{ch20_exchange_1921}{Operators at the Garfield exchange in Seattle,
about 1921. Each position served a group of lines; how many positions and trunks an exchange needed for a
given grade of service was exactly the question Erlang answered.}

\begin{history}[Agner Krarup Erlang]
A.~K.~Erlang (1878--1929) joined the Copenhagen Telephone Company in 1908 as its first scientific collaborator, at
the invitation of its chief engineer, who wanted to know how many lines and switches the growing exchanges needed.
In 1909 Erlang showed that telephone call arrivals follow a Poisson distribution, and in 1917 he published the
loss formula that bears his name, together with the formula for systems in which calls wait instead of being lost.
He worked alone and published mostly in Danish, but his results spread through the British Post Office and Bell
System and became the foundation of queueing theory and teletraffic engineering. The international unit of
traffic intensity was named the \emph{erlang} in his honour in 1946.
\end{history}

\begin{analogy}[How many checkout lanes?]
A supermarket manager must decide how many checkout lanes to open. Too few, and at the busy hour shoppers find
every lane jammed and walk out; too many, and cashiers sit idle. Customers arrive at random, and each takes a
random time to serve. Erlang B answers the version in which a customer who finds every lane busy simply leaves
(a blocked call is lost); Erlang C answers the version in which they queue. The analogy's limit: shoppers can
see the queues and pick the shortest, and some come back later, whereas Erlang's callers are assumed to
arrive independently of how busy the system is.
\end{analogy}

\mdcphoto[0.55]{ch20_checkout}{Checkout lanes in a supermarket. How many to staff at the busy hour, so that
few customers find every lane occupied, is the same mathematical question as how many radio channels a cell
needs.}

\subsection{Offered traffic and the erlang}

Calls arrive at a cell as a Poisson process of rate $\lambda$ calls per unit time, and each lasts a random
\term{holding time} with mean $h$. The \term{offered traffic} is
\begin{equation}
A=\lambda h\quad\text{erlangs}.
\end{equation}
One erlang is one channel continuously occupied; equivalently, $A$ is the average number of calls that would be in
progress if no call were ever refused. A subscriber who makes one call of two minutes in the \term{busy hour}, the
hour of the day with the highest traffic, generates $1\times2/60=0.033$\,E, or 33\,mE. The \term{grade of service}
(GoS) is the probability that a call attempt in the busy hour is blocked; 1--2\% is the traditional target for
voice.

\subsection{Erlang B: blocked calls cleared}

If blocked calls simply disappear (the caller gives up, or tries again much later), a cell with $C$ channels is the
queue M/M/$C$/$C$: Poisson arrivals, exponential holding times (for now), $C$ servers, no waiting room. Its state is
the number $k$ of calls in progress, a birth--death process (Figure~\ref{fig:ch20_birthdeath}) with birth rate
$\lambda$ in every state $k<C$ and death rate $k/h$ in state $k$.

\begin{figure}[htbp]\centering
\begin{tikzpicture}[font=\sffamily\scriptsize,>={Stealth[length=2mm]},
  st/.style={circle,draw=navy,fill=navy!8,thick,minimum size=0.75cm,inner sep=0pt}]
\foreach \k/\x in {0/0,1/1.9,2/3.8} {\node[st] (s\k) at (\x,0) {\k};}
\node (dots) at (5.5,0) {$\cdots$};
\node[st] (sc1) at (7.2,0) {$C\!-\!1$};
\node[st,draw=accent,fill=accent!10] (sc) at (9.1,0) {$C$};
\foreach \a/\b/\m in {s0/s1/1,s1/s2/2,sc1/sc/C} {
  \draw[->,thick,navy] (\a) to[bend left=35] node[above]{$\lambda$} (\b);
  \draw[->,thick,green!50!black] (\b) to[bend left=35] node[below]{$\m/h$} (\a);}
\draw[->,thick,navy] (s2) to[bend left=35] node[above]{$\lambda$} (dots);
\draw[->,thick,green!50!black] (dots) to[bend left=35] node[below]{$3/h$} (s2);
\draw[->,thick,navy] (dots) to[bend left=35] node[above]{$\lambda$} (sc1);
\draw[->,thick,green!50!black] (sc1) to[bend left=35] node[below]{$(C\!-\!1)/h$} (dots);
\draw[->,thick,accent] (sc.east) -- ++(1.0,0) node[right,align=left]{arrivals\\blocked};
\node[text=gray] at (4.55,-1.25) {state $k$ = calls in progress; each call ends at rate $1/h$, so $k$ calls end at rate $k/h$};
\end{tikzpicture}
\caption{The Erlang loss system as a birth--death chain. Calls arrive at rate $\lambda$ whenever a channel is
free; with $k$ calls in progress, one of them ends at rate $k/h$. In state $C$ every channel is busy and new
calls are lost.}
\label{fig:ch20_birthdeath}
\end{figure}

In equilibrium the probability flow across each cut balances,
$\lambda\,P_{k-1}=(k/h)\,P_k$, so $P_k=P_{k-1}A/k$ and
\begin{equation}
P_k=\frac{A^k/k!}{\sum_{j=0}^{C}A^j/j!},\qquad k=0,1,\dots,C,
\end{equation}
a Poisson distribution truncated at $C$. Because Poisson arrivals see time averages, the probability that an arriving
call is blocked is the probability that all channels are busy:
\begin{equation}
B(A,C)=\frac{A^C/C!}{\sum_{k=0}^{C}A^k/k!}.
\label{eq:ch20:erlangb}
\end{equation}
This is the \term{Erlang B formula}. Figure~\ref{fig:ch20_switchboard} watches it at work on a simulated
switchboard. A remarkable property makes it robust: it is \emph{insensitive} to the
distribution of the holding times, depending only on their mean, so it holds for real call durations, which are far
from exponential. For computation, factorials overflow quickly; the recursion
\begin{equation}
B(A,0)=1,\qquad B(A,c)=\frac{A\,B(A,c-1)}{c+A\,B(A,c-1)}
\label{eq:ch20:erlangbrec}
\end{equation}
is numerically stable and is how every Erlang calculator works. The \term{carried traffic} is $A(1-B)$, and the
channel utilisation is $A(1-B)/C$.

\mdcfig{ch20_switchboard}{A live switchboard. (a) Four simulated hours of a 10-line trunk offered 7\,E of
three-minute calls: the number of lines in use wanders, and each time it touches the ceiling, arriving callers
(crosses) are turned away. (b) The long-run distribution of lines in use is Erlang's truncated Poisson law; the
red bar, the probability that all lines are busy, is the blocking probability $B(7,10)\approx7.9\%$. The short
simulation turned away about 10\% of its callers, within the run-to-run scatter of a four-hour sample.}

\mdcfig{ch20_erlang}{Traffic theory. (a) Erlang B blocking probability versus offered traffic for $C=1$ to 100
channels. (b) Trunking efficiency: the carried traffic per channel at a fixed grade of service rises steeply with the
number of channels; large trunks are efficient. (c) Erlang C probability that a call must wait, versus utilisation
$A/C$: with many servers, waiting stays rare until the utilisation is high.}

\subsection{Trunking efficiency}

Table~\ref{tab:ch20:erlangb} and Figure~\ref{fig:ch20_erlang}b show the most important property of the
formula: \term{trunking efficiency}, the economies of scale of shared resources. At 2\% blocking, 10 channels carry
5.1\,E, a utilisation of about 50\%, while 100 channels carry 88\,E, about 86\% (Figure~\ref{fig:ch20_trunk_pools}).
Ten separate 10-channel cells carry 51\,E between them; one 100-channel cell carries 88\,E. Random fluctuations of the
number of active calls grow only as the square root of the mean, so a large group needs proportionally less headroom.
Every technique that splits a pool of channels into smaller pools---sectorization, separate carriers for separate
services, dedicated channels for particular users---pays a trunking penalty, and every technique that pools resources
across cells, carriers or technologies gains. A single snake queue feeding ten checkouts beats ten separate queues for
the same reason.

\mdcfig[0.5]{ch20_trunk_pools}{Trunking efficiency in one picture: the same 100 lines carry 51\,E at 2\%
blocking when split into ten exchanges, but 88\,E when pooled.}

\begin{worked}[Dimensioning an AMPS cell, and the price of sectors]
An AMPS cell with $N=7$ has 57 voice channels. Each subscriber generates 0.033\,E in the busy hour (one two-minute
call). How many subscribers can the cell serve at 2\% blocking (a) as one omnidirectional cell, (b) as three
sectors of 19 channels each, and (c) with three sectors and $N=4$?

(a) Solving $B(A,57)=0.02$ gives $A=46.8$\,E, so the cell serves $46.8/0.033\approx1400$ subscribers.

(b) Each sector is now a separate trunk of 19 channels, carrying 12.3\,E at 2\%: $3\times12.3=37.0$\,E in all,
about 1110 subscribers. Sectoring for interference costs 21\% of the cell's capacity in trunking efficiency.

(c) The sectors' interference advantage allows $N=4$ instead of 7 (Section~\ref{sec:ch20:cellular}), so each cell
gets about $395/4\approx99$ voice channels, 33 per sector. Each sector carries 24.6\,E, the cell 73.9\,E, about 2200
subscribers, 58\% more than the omnidirectional $N=7$ cell. The lesson is typical of cellular engineering:
sectorization pays for itself only when its SIR gain is spent on tighter reuse.
\end{worked}

\begin{tryit}[Run a switchboard]
Lab 26's \emph{Erlang: a live switchboard} animates calls arriving and leaving. Start at 57 \emph{Channels} and
40\,E of \emph{Offered traffic}, then push the traffic up until the blocking counter reaches the 2\%
\emph{Grade of service}. Now set \emph{Split the channels into sectors} to 3 and watch the blocking jump at the
same total load: that is case (b) of the worked example. In Lab 11's \emph{Trunking: a live switchboard},
\emph{Split into separate groups} does the same for a smaller trunk.
\end{tryit}

\begin{table}[htbp]\centering\small
\caption{Erlang B: offered traffic (erlangs) that $C$ channels carry at a given blocking probability.}
\label{tab:ch20:erlangb}
\begin{tabular}{@{}r rrrr @{\hspace{2.5em}} r rrrr@{}}
\toprule
$C$ & 0.5\% & 1\% & 2\% & 5\% & $C$ & 0.5\% & 1\% & 2\% & 5\%\\
\midrule
1 & 0.005 & 0.010 & 0.020 & 0.053 & 15 & 7.38 & 8.11 & 9.01 & 10.6\\
2 & 0.105 & 0.153 & 0.223 & 0.381 & 20 & 11.1 & 12.0 & 13.2 & 15.2\\
3 & 0.349 & 0.455 & 0.602 & 0.899 & 25 & 15.0 & 16.1 & 17.5 & 20.0\\
4 & 0.701 & 0.869 & 1.09 & 1.52 & 30 & 19.0 & 20.3 & 21.9 & 24.8\\
5 & 1.13 & 1.36 & 1.66 & 2.22 & 40 & 27.4 & 29.0 & 31.0 & 34.6\\
6 & 1.62 & 1.91 & 2.28 & 2.96 & 50 & 36.0 & 37.9 & 40.3 & 44.5\\
7 & 2.16 & 2.50 & 2.94 & 3.74 & 60 & 44.8 & 47.0 & 49.6 & 54.6\\
8 & 2.73 & 3.13 & 3.63 & 4.54 & 75 & 58.2 & 60.7 & 63.9 & 69.7\\
10 & 3.96 & 4.46 & 5.08 & 6.22 & 100 & 80.9 & 84.1 & 88.0 & 95.2\\
12 & 5.28 & 5.88 & 6.61 & 7.95 & 150 & 127.4 & 131.6 & 136.8 & 146.7\\
\bottomrule
\end{tabular}
\end{table}

\subsection{Erlang C: blocked calls delayed}

If a call that finds all channels busy waits in a queue instead (as in a call centre, or for a packet waiting for
a scheduler), the system is the M/M/$C$ queue. The balance equations are the same for $k\le C$, and for $k>C$ the
death rate stays at $C/h$, so $P_k=P_C\,(A/C)^{k-C}$. The queue is stable only if $A<C$. Summing the geometric
tail gives the probability that an arriving call must wait, the \term{Erlang C formula}:
\begin{equation}
C(A,C)=\frac{\dfrac{A^C}{C!}\,\dfrac{C}{C-A}}{\displaystyle\sum_{k=0}^{C-1}\frac{A^k}{k!}+\frac{A^C}{C!}\,\frac{C}{C-A}}
=\frac{C\,B(A,C)}{C-A\,\big(1-B(A,C)\big)} .
\label{eq:ch20:erlangc}
\end{equation}
With first-come first-served queueing and exponential holding times, the waiting time $W$ satisfies
\begin{equation}
\Pr\{W>t\}=C(A,C)\,e^{-(C-A)t/h},\qquad \E[W]=\frac{C(A,C)\,h}{C-A}.
\end{equation}
Unlike Erlang B, Erlang C does depend on the holding-time distribution, so it is an approximation for real traffic.
Figure~\ref{fig:ch20_erlang}c shows the trunking effect again: a single server at 80\% utilisation makes 80\% of
callers wait, whereas 50 servers at 80\% utilisation make only a few percent wait. With 30 agents and 25\,E of
offered traffic (83\% utilisation) and three-minute calls, 25\% of callers wait and the mean wait over all callers
is about 9\,s (Figure~\ref{fig:ch20_erlang_wait}). Two agents fewer and the wait more than doubles; three more
and it nearly vanishes---the steepness that makes call-centre staffing, and scheduler dimensioning, so sensitive.

\chTwentyFigH[0.6]{ch20_erlang_wait}{Erlang C for the call-centre example: the fraction of callers who wait
longer than $t$, with 25\,E of offered traffic and three-minute calls, for three staffing levels. Near full
utilisation, every agent matters.}

\begin{inpractice}[Erlang in the packet era]
Voice over LTE and 5G no longer uses dedicated circuits, so is Erlang obsolete? Not quite. Operators still
dimension VoLTE capacity, signalling channels, paging capacity and RACH resources in erlang-like terms, because each
of these is a shared pool serving random requests. The insensitivity of Erlang B makes it a good first model for any
resource that is held for a while and then released---licences, radio bearers, hardware channel elements, even GPUs
in a server farm. For finite populations, where each busy user reduces the arrival rate, the Engset formula replaces
Erlang B; for retrial behaviour, more elaborate models exist. In packet schedulers, the corresponding question is
not ``blocked or not'' but ``how much throughput does each user get'', which leads to the processor-sharing models
used in modern capacity planning: with $n$ active users sharing a cell of capacity $c$ fairly, a user's throughput is
about $c/n$, and the number of active users follows the same Poisson-type statistics as calls in Erlang's exchange.
\end{inpractice}

%======================================================================
\section{Interference management}
\label{sec:ch20:interference}
%======================================================================

Fixed reuse patterns trade spectrum for interference once and for all, like neighbours who agree never to mow
their lawns on the same day of the week. Modern systems run at reuse 1---everyone mows whenever they like---and
manage interference dynamically, with power control, coordination between neighbouring cells, and, in the limit,
joint transmission from several sites.

\mdcfig[0.5]{ch20_fpc}{Fractional power control: transmit power versus path loss for three values of $\alpha$
(with $P_0$ chosen so the curves meet at 100\,dB). Full compensation ($\alpha=1$) makes far users shout
loudest; $\alpha<1$ lets them accept a lower received level and spares the neighbouring cells.}

\subsection{Power control}

On the uplink, terminals near the base station would be received tens of decibels stronger than those at the edge
if all transmitted at full power. In a CDMA system this near--far effect is fatal and power control must equalise
received powers to within about a decibel, at rates up to 1500 updates per second in UMTS
(Chapter~\ref{ch:spreadspectrum}). In OFDMA, users in the same cell are orthogonal, so near--far within a cell
matters less (it still matters for the dynamic range of the receiver and for leakage between adjacent resource
blocks), but every terminal's power is interference in the \emph{neighbouring} cells. LTE and NR use
\term{fractional power control}. The transmit power of the uplink shared channel in a slot is, in simplified form
(TS~36.213 for LTE, TS~38.213 for NR),
\begin{equation}
P=\min\Big\{P_{\max},\;P_0+10\log_{10}M+\alpha\cdot\mathrm{PL}+\Delta_{\mathrm{TF}}+f\Big\}\quad\text{dBm},
\label{eq:ch20:fpc}
\end{equation}
where $M$ is the number of allocated resource blocks, PL the terminal's estimate of its downlink path loss,
$\Delta_{\mathrm{TF}}$ an adjustment for the modulation and coding scheme, and $f$ an accumulated closed-loop
correction driven by transmit-power-control commands.

The \emph{open-loop} part, $P_0+\alpha\,\mathrm{PL}$, is the interesting one (Figure~\ref{fig:ch20_fpc}). With
$\alpha=1$ (full compensation) every user arrives with the same power spectral density $P_0$, as in CDMA, so the
cell-edge users, who have the highest path loss, transmit the most power and create the most interference in
neighbouring cells. With $\alpha<1$ (typical values 0.7--0.8), received power at the serving cell grows with
proximity, so near users get higher SINR and rates, edge users accept lower SINR, and the interference generated in
neighbouring cells drops. Fractional control trades a little cell-edge throughput for a larger gain in average
throughput, and it is one of the simplest and most effective tools for reuse-1 uplinks. Downlink power control is less
important in OFDMA: the base station's power is shared among resource blocks, and link adaptation, not power,
matches each user's rate to its channel.

\mdcfig{ch20_ffr}{Fractional frequency reuse. (a) Strict FFR: a common band at reuse 1 for centre users plus three
edge sub-bands at reuse 3. (b) Soft frequency reuse: every cell uses all sub-bands, with full power on its own edge
sub-band. (c)~SINR and (d) per-user rate distributions in a hexagonal network ($n=3.5$, 6\,dB shadowing, 20\,dB
cell-edge SNR). The 30\% of users with the worst reuse-1 SINR are classed as edge users; in (d) each class shares its
part of the band equally. FFR raises the fifth-percentile rate about fourfold relative to reuse 1 at a cost of about
13\% in mean rate; reuse 3 everywhere helps the edge less and costs a third of the mean.}

\subsection{Inter-cell interference coordination}

\term{Inter-cell interference coordination} (ICIC) lets neighbouring cells agree on which parts of the band each
will use for its vulnerable cell-edge users. The simplest static scheme is \term{fractional frequency reuse} (FFR):
split the band into a common part used by every cell for users near the centre, at reuse 1, and an edge part
divided into three sub-bands used at reuse 3 by the cell-edge users (Figure~\ref{fig:ch20_ffr}a). Each cell-edge
user then sees interference only from a third of the neighbouring cells. In \term{soft frequency reuse} (SFR), every
cell uses the whole band, but transmits at full power only on its own third and at reduced power elsewhere, with
the edge users scheduled on the high-power third (Figure~\ref{fig:ch20_ffr}b). LTE Release~8 standardised the
messages that support dynamic versions of these schemes over the X2 interface between base stations: for the
downlink, a cell announces on which resource blocks it intends to transmit at high power (relative narrowband
transmit power, RNTP); for the uplink, it warns its neighbours which resource blocks its edge users will occupy
(high interference indicator) and complains about those on which it suffers (overload indicator).

Figure~\ref{fig:ch20_ffr} quantifies the trade-off in a simple model. Pure reuse 1 gives the best mean rate but leaves
the worst 5\% of users with a rate below 0.1\,b/s/Hz. Reuse 3 everywhere lifts the edge but wastes two-thirds of the
band at the centre. FFR does both: the edge users' SINR moves from the reuse-1 curve to the reuse-3 curve, and the
centre users keep the full-band efficiency. The best split between the common and edge bands depends on the user
distribution and the traffic, which is why later systems moved from static plans to coordinated scheduling.

\mdcfig{ch20_range_expansion}{Cell range expansion along a line through a macro site (46\,dBm) and a small cell
(30\,dBm) 700\,m away, using 3GPP-style macro and pico path-loss laws. Users served by the strongest
downlink join the small cell only inside its natural area (green); a 9\,dB association bias extends it
(orange), offloading users who then rely on almost blank subframes to escape the stronger macro signal.}

\subsection{Heterogeneous networks: range expansion and almost-blank subframes}

A small cell under a macro cell faces a mismatch. Users connect to the cell with the strongest downlink, and the
macro cell, transmitting perhaps 16\,dB more power, wins over most of the small cell's natural area. But the
uplink path loss, which is what matters for the terminal's battery and for interference, is lowest to the
\emph{nearest} cell, often the small one. LTE Release~10 introduced \term{enhanced ICIC} (eICIC) for this case.
\term{Cell range expansion} adds a bias, typically a few decibels and up to about 9\,dB, to the small cell's
measured power when choosing the serving cell, so more users offload onto it (Figure~\ref{fig:ch20_range_expansion}).
The users in the expanded region then receive the macro cell's signal more strongly than their own, so the macro
cell transmits \term{almost blank subframes} (ABS), subframes in which it sends only essential reference signals
and no data, and the small cell schedules its range-expanded users in those subframes---the loud neighbour agrees
to keep quiet for a few minutes every hour so that the quiet one can be heard. Further enhancements (FeICIC,
Release~11) added receiver cancellation of the residual macro reference signals. NR, with its leaner always-on
signals, has much less residual interference to cancel, and relies on scheduling and beamforming.

\subsection{Coordinated multipoint}

The most ambitious form of interference management treats neighbouring sites as one distributed antenna array.
In \term{coordinated multipoint} (CoMP, LTE Release~11; multi-TRP transmission in NR), several transmission
points may \emph{jointly transmit} the same data to a cell-edge user (joint transmission, turning interference into
signal); \emph{dynamically select} which point serves the user from one slot to the next; or
\emph{coordinate their scheduling and beamforming} so that each one's beams avoid the others' users
(Figure~\ref{fig:ch20_comp}). The gains of joint transmission can be large in theory, since it removes inter-cell
interference altogether, but they depend on accurate channel knowledge from several sites and on fast,
high-capacity backhaul between them; field gains have been modest except where the points share a baseband unit.
Its logical conclusion, \emph{cell-free massive MIMO}, in which many distributed antennas jointly serve all users
with no cells at all, is discussed in Chapter~\ref{ch:mimo}.

\begin{figure}[htbp]\centering
\begin{tikzpicture}[font=\sffamily\scriptsize,>={Stealth[length=2mm]},
  bs/.style={regular polygon,regular polygon sides=3,draw=navy,fill=navy!60,minimum size=0.45cm,inner sep=0pt},
  ue/.style={circle,draw=accent,fill=accent!60,minimum size=0.25cm,inner sep=0pt}]
\begin{scope}
\node[font=\sffamily\small\bfseries,text=navy] at (1.5,1.55) {joint transmission};
\node[bs] (a) at (0,0) {}; \node[bs] (b) at (3,0) {}; \node[ue] (u) at (1.5,-0.9) {};
\draw[->,thick,green!50!black] (a) -- (u); \draw[->,thick,green!50!black] (b) -- (u);
\node[align=center,text=gray] at (1.5,-1.65) {both sites send the same data:\\interference becomes signal};
\end{scope}
\begin{scope}[xshift=4.6cm]
\node[font=\sffamily\small\bfseries,text=navy] at (1.5,1.55) {dynamic point selection};
\node[bs] (a) at (0,0) {}; \node[bs] (b) at (3,0) {}; \node[ue] (u) at (1.5,-0.9) {};
\draw[->,thick,green!50!black] (b) -- (u); \draw[->,thick,gray,dashed] (a) -- (u);
\node[align=center,text=gray] at (1.5,-1.65) {the better site serves this slot;\\the other stays silent on it};
\end{scope}
\begin{scope}[xshift=9.2cm]
\node[font=\sffamily\small\bfseries,text=navy] at (1.5,1.55) {coordinated beamforming};
\node[bs] (a) at (0,0) {}; \node[bs] (b) at (3,0) {}; \node[ue] (u) at (0.5,-1.0) {}; \node[ue] (v) at (2.5,-1.0) {};
\fill[green!50!black,opacity=0.25] (a.south) -- (0.15,-1.35) -- (0.85,-1.35) -- cycle;
\fill[green!50!black,opacity=0.25] (b.south) -- (2.15,-1.35) -- (2.85,-1.35) -- cycle;
\node[align=center,text=gray] at (1.5,-1.75) {each site steers its beam\\away from the other's user};
\end{scope}
\end{tikzpicture}
\caption{Three flavours of coordinated multipoint for a user at the boundary between two sites. They differ in
how much the sites must share: user data and channel knowledge (joint transmission), fast scheduling decisions
(dynamic point selection), or only channel knowledge and beam choices (coordinated beamforming).}
\label{fig:ch20_comp}
\end{figure}

%======================================================================
\section{The capacity of cellular networks}
\label{sec:ch20:capacity}
%======================================================================

How much data can a cellular network carry per square kilometre, and how does that change as base stations are
added? Hexagonal models answer such questions by simulation, one geometry at a time. A different approach, which
has reshaped network theory since about 2010, models the base-station locations as random---and, surprisingly,
the randomness is what makes the problem solvable with pencil and paper.

\subsection{Base stations as a Poisson point process}

Real base stations are not on a lattice. They sit where an operator could obtain a site, and their density follows
the population. Andrews, Baccelli and Ganti proposed in 2011 to model them as a homogeneous \term{Poisson point
process} (PPP) $\Phi$ of density $\lambda$ base stations per unit area: the number in any region of area $A$ is Poisson
with mean $\lambda A$, independently for disjoint regions (Figure~\ref{fig:ch20_ppp}a). Picture throwing a handful of
rice onto a map: each grain is a base station, and each point of the map is served by the nearest grain. The model
sounds crude, but it captures the irregularity of real networks and, crucially, it makes the interference tractable.

Place a typical user at the origin, served by its nearest base station at distance $r$. With Rayleigh fading
($h\sim$ exponential with unit mean on every link), path loss $d^{-\alpha}$, equal transmit power $P$ and noise power
$\sigma^2$, the downlink SINR is
\begin{equation}
\mathrm{SINR}=\frac{h\,r^{-\alpha}}{I+\sigma^2/P},\qquad
I=\sum_{x\in\Phi\setminus\{b_0\}}g_x\,\|x\|^{-\alpha}.
\end{equation}
\mdcfig{ch20_ppp}{Stochastic geometry of a cellular network. (a) Base stations drawn from a Poisson point process,
with the Voronoi cells of nearest-station association. (b) Coverage probability with Rayleigh fading and no noise:
the closed-form PPP result~\eqref{eq:ch20:pppnoisefree} for three path-loss exponents, a Monte Carlo check for
$\alpha=4$, and a hexagonal grid of the same density, which is about 4\,dB more optimistic at the median. (c) Coverage
versus density when noise matters ($\alpha=4$, 10\,dB SNR at 300\,m): densification moves the network from noise-limited
to interference-limited operation, after which coverage no longer depends on density (dotted lines).}

The \term{coverage probability} $P_c(T)=\Pr\{\mathrm{SINR}>T\}$ is the complementary CDF of SINR. Averaging over
the distance to the serving station and over the random interferers (box below) gives the general result
\begin{equation}
P_c(T)=\pi\lambda\int_0^\infty e^{-\pi\lambda v\,(1+\rho(T,\alpha))-T\sigma^2v^{\alpha/2}/P}\,dv,
\label{eq:ch20:ppp}
\end{equation}
and, when interference dominates noise,
\begin{equation}
P_c(T)=\frac{1}{1+\rho(T,\alpha)},\qquad
\rho(T,4)=\sqrt T\left(\frac\pi2-\arctan\frac{1}{\sqrt T}\right).
\label{eq:ch20:pppnoisefree}
\end{equation}
For $\alpha=4$ and $T=1$ (0\,dB), $P_c=1/(1+\pi/4)\approx0.56$: in a reuse-1 network, about 44\% of locations have an
SINR below 0\,dB. Integrating the coverage curve gives the mean spectral efficiency, $\E[\ln(1+\mathrm{SINR})]\approx1.49$
nats/s/Hz, or about 2.15\,b/s/Hz, for $\alpha=4$.

\begin{deeper}[Three facts that make the PPP computable]
\begin{enumerate}
\item The distance to the nearest point of a PPP has density $f(r)=2\pi\lambda r\,e^{-\pi\lambda r^2}$, since the
probability of no point within $r$ is $e^{-\lambda\pi r^2}$.
\item Rayleigh fading turns the coverage condition into an exponential tail:
$\Pr\{h>Tr^{\alpha}(I+\sigma^2/P)\}=\E\big[e^{-Tr^{\alpha}(I+\sigma^2/P)}\big]=e^{-Tr^\alpha\sigma^2/P}\,\mathcal L_I(Tr^\alpha)$,
where $\mathcal L_I(s)=\E[e^{-sI}]$ is the Laplace transform of the interference.
\item The Laplace transform of a sum over a PPP has a closed form (the probability generating functional). For the
interferers, which lie outside the disk of radius $r$,
\begin{equation}
\mathcal L_I(s)=\exp\!\left(-2\pi\lambda\int_r^\infty\Big(1-\frac{1}{1+s\,v^{-\alpha}}\Big)v\,dv\right)
=\exp\!\big(-\pi\lambda r^2\rho(T,\alpha)\big)\quad\text{at } s=Tr^\alpha,
\end{equation}
with $\rho(T,\alpha)=T^{2/\alpha}\int_{T^{-2/\alpha}}^{\infty}\frac{du}{1+u^{\alpha/2}}$ after the substitution
$u=(v/r)^2T^{-2/\alpha}$.
\end{enumerate}
Multiplying the three pieces and integrating over $r$ (with $v=r^2$) gives~\eqref{eq:ch20:ppp}; setting $\sigma^2=0$
makes the integral elementary and gives~\eqref{eq:ch20:pppnoisefree}.
\end{deeper}

\chTwentyFigH{ch20_densify_maps}{Scale invariance, seen. (a) Downlink SIR over an $8\times8$\,km patch of a PPP network
($\alpha=4$, no fading). (b) A network four times denser, shown over a $4\times4$\,km patch: shrink the map by two and
it is statistically the same picture. The median SIR is about 2.5\,dB in both, yet (b) serves four times as many
cells' worth of traffic per square kilometre.}

\begin{keyidea}[In an interference-limited network, density does not change the SINR]
Equation~\eqref{eq:ch20:pppnoisefree} does not contain $\lambda$. Adding base stations brings the serving station
closer, but it brings the interferers closer by exactly the same factor, and the SINR distribution is unchanged
(Figure~\ref{fig:ch20_densify_maps}). Each new base station therefore adds a full cell's worth of capacity: the
\term{area spectral efficiency},
\begin{equation}
\mathrm{ASE}=\lambda\cdot\E[\log_2(1+\mathrm{SINR})]\quad\text{b/s/Hz/km}^2,
\end{equation}
grows linearly with density. With 10 base stations per square kilometre and $\alpha=4$, the ASE is about
$10\times2.15=21.5$\,b/s/Hz/km$^2$, or 2.15\,Gb/s/km$^2$ with 100\,MHz of spectrum; ten times the density gives ten
times the capacity. This scale invariance is the formal statement of the cellular concept's promise.
\end{keyidea}

\begin{tryit}[Densify a random network]
Lab 26's \emph{Stochastic geometry: PPP coverage} draws a random deployment and its SINR map. With \emph{SNR at 300\,m}
set high (30\,dB or more), drag \emph{Base-station density} across two decades: the map shrinks but the coverage curve
does not move, and the challenge value $1/(1+\pi/4)$ appears at $T=0$\,dB. Now lower the SNR to 0\,dB and repeat:
at low density, adding base stations helps, until the curve meets the interference-limited bound. Press
\emph{New deployment} to see how much one random draw varies.
\end{tryit}

\chTwentyFigH[0.62]{ch20_dual_slope}{Densification with one and with two path-loss slopes (Monte Carlo, PPP
base stations, Rayleigh fading, no noise). With a single exponent the area spectral efficiency grows in
proportion to density, as the key idea promised. When links shorter than 50\,m follow the gentler free-space
law, interferers start to arrive line of sight once cells shrink below that size, and the growth slows: at
10\,000 sites per km$^2$ the dual-slope network carries about a quarter of the single-slope figure.}

\subsection{Where densification stops paying}

The scale invariance rests on assumptions that fail at very high density, and the failures are instructive.
\begin{itemize}
\item \textbf{Noise at low density.} When cells are large, the cell edge is noise-limited and adding base stations
improves coverage (Figure~\ref{fig:ch20_ppp}c). This is the coverage-driven phase of network roll-out.
\item \textbf{Path loss is not a single power law.} At short range, links are increasingly line of sight with a
path-loss exponent near 2, while the interferers further away remain obstructed. Zhang and Andrews showed in 2015 that
with such multi-slope path loss, the SINR \emph{falls} as the network densifies past the point where interferers enter
the low-exponent region, and the ASE growth slows or stops. The physical cure is to keep interferers obstructed, by
lower antennas, building penetration loss or narrow beams.
\item \textbf{Antenna heights.} Even at zero horizontal distance, a user is a few metres below the antenna, so path
loss is bounded; at very high density the serving signal stops improving while interference keeps growing.
\item \textbf{Load.} If base stations outnumber active users, many cells are empty. An idle base station that switches
off its transmitter creates no interference, which makes ultra-dense networks \emph{better} than the full-load model
predicts, provided the always-on signals are lean (one of NR's design goals).
\item \textbf{Cost.} Each base station needs a site, power and backhaul, and each user in a tiny cell hands over
constantly. In practice the economics of sites and fibre, not the SINR, limits densification.
\end{itemize}

\begin{inpractice}[Hexagons, Poisson points and real networks]
The hexagonal model is an optimistic bound: interferers are as far away as they could be. The PPP is a pessimistic
one: base stations can be arbitrarily close together. Comparisons with real deployments (Andrews, Baccelli and Ganti
used an urban macro network in their 2011 paper) place measured SINR distributions between the two, often closer to
the PPP, and a constant SINR offset of a couple of decibels maps one model onto the other surprisingly well. Industry
system-level simulators, following 3GPP's evaluation methodology (for example TR~36.814 and TR~38.901), still use
hexagonal layouts of 19 or 57 sectors with wrap-around, because standardisation needs a common, reproducible
geometry. Theorists use the PPP because it gives formulas. A good engineer uses both and checks each against the
other.
\end{inpractice}

%======================================================================
\section{Mobility: handover and location management}
\label{sec:ch20:mobility}
%======================================================================

Drive across a city while on a video call and your phone may change cells dozens of times without a glitch you
can notice. Behind that seamlessness sits a careful compromise between switching too eagerly and switching too
late. A cellular network must keep track of terminals that move, both while they are in a call (handover) and
while they are idle (location management and paging).

\mdcfig{ch20_a3}{Event A3 in slow motion. The neighbour's filtered RSRP crosses the serving cell's long before
the handover: first it must exceed the serving level plus the hysteresis, and then stay there for the
time-to-trigger (shaded). Only then is the measurement report sent and the handover prepared.}

\subsection{Handover}

When a connected terminal moves from one cell to another, the connection must be transferred, ideally without the
user noticing. In a \term{hard handover} (``break before make''), the terminal leaves the old cell and then joins the
new one, with a brief interruption: GSM, LTE and NR use hard handover, with interruptions of a few tens of
milliseconds that voice codecs and TCP tolerate. In a \term{soft handover} (``make before break''), used by CDMA
systems, the terminal communicates with several base stations at once during the transition, and their signals
are combined; it gains macrodiversity but uses resources in several cells (Chapter~\ref{ch:spreadspectrum}). NR
Release~16 added a make-before-break option for latency-critical services, \emph{dual active protocol stack}
handover, in which the terminal keeps its old link until the new one carries data.

The decision is based on measurements. Terminals measure the reference signals of their serving cell and of
neighbours (in LTE and NR, the reference signal received power, RSRP, and quality, RSRQ), average them to remove
fast fading, and filter them further with a first-order recursion configured by the network,
\begin{equation}
F_n=(1-a)\,F_{n-1}+a\,M_n,\qquad a=2^{-k/4},
\end{equation}
before evaluating reporting criteria. The workhorse criterion is LTE and NR \emph{event A3}: a neighbour becomes
better than the serving cell by an offset. Stripped of its cell-specific offsets, it fires when
\begin{equation}
M_{\text{neighbour}}>M_{\text{serving}}+\mathrm{Hys}\quad\text{continuously for a time-to-trigger } T_{\text{TTT}}.
\end{equation}
The terminal then reports, the serving cell prepares the target cell over the inter-base-station interface, and
commands the handover, after which the terminal performs contention-free random access in the target cell
(Figure~\ref{fig:ch20_a3}).

\subsection{Hysteresis, time-to-trigger and the ping-pong problem}

\mdcfig{ch20_handover}{Handover between two sites 2\,km apart (path-loss exponent 3.5, 8\,dB shadowing with 50\,m
decorrelation distance and inter-site correlation 0.5, L3 filter with $a=0.5$, user speed 15\,m/s). (a) Filtered
received power from the two cells along one drive, with the handovers triggered by a zero-hysteresis rule (red ticks)
and by 3\,dB hysteresis with a 1.28\,s time-to-trigger (orange triangles). (b) Mean number of handovers per drive versus
hysteresis for three time-to-trigger values. (c) Fraction of the route on which the serving cell is more than 3\,dB
weaker than the best cell: the price of suppressing ping-pong.}

\begin{analogy}[Not switching queues every time the other one looks shorter]
At the supermarket the next queue always seems to move faster. Shoppers who hop to it every time it looks
shorter spend their afternoon hopping, and often end up behind where they started. Sensible shoppers switch only
when the other queue is clearly shorter (a hysteresis margin) and stays that way for a while (a time-to-trigger).
The cost of their caution is occasionally staying a little too long in a slow queue. A phone near a cell boundary
faces exactly this choice, many times a minute.
\end{analogy}

Shadowing makes the measured powers of two cells cross many times near the boundary
(Figure~\ref{fig:ch20_handover}a). A rule that always picks the strongest cell would hand over back and forth at every
crossing, a \term{ping-pong} that wastes signalling, interrupts the user and risks failures. The two classical
defences are a \term{hysteresis} margin, which requires the new cell to be better by several decibels, and a
\term{time-to-trigger}, which requires the condition to persist. Both delay the handover, and a late handover keeps the
user on a cell that is getting worse, eventually causing a radio link failure. Figure~\ref{fig:ch20_handover}b--c
quantifies the trade-off for a user driving at 15\,m/s between two sites with 8\,dB of correlated shadowing: with no
hysteresis and no time-to-trigger there are dozens of handovers per drive; 3--4\,dB of hysteresis with a few hundred
milliseconds of time-to-trigger brings this to one or two, at the price of a few percent of the route spent on a cell
more than 3\,dB weaker than the best. Typical network settings are of this order, with the standard allowing hysteresis
up to 15\,dB and time-to-trigger values from 0 to 5.12\,s.

\begin{tryit}[Tune a handover]
Open Lab 26's \emph{Handover: hysteresis and TTT}. With \emph{Hysteresis} and \emph{Time-to-trigger} at zero, press
\emph{Another drive} a few times and count the ping-pongs. Raise \emph{Hysteresis} to 3\,dB, then add 320\,ms of
\emph{Time-to-trigger}; the handovers collapse to one or two. Now raise \emph{Speed} to 40\,m/s: the same
time-to-trigger now covers three times the distance, and the late-handover penalty grows. Networks tune these
settings per speed class for exactly this reason.
\end{tryit}

Because the right settings depend on speed and environment, networks tune them automatically. \emph{Mobility robustness
optimisation}, one of the self-organising-network functions standardised by 3GPP, counts handovers that were too late
(failure before or during handover), too early (failure soon after) or to the wrong cell, and adjusts offsets per
neighbour relation. NR Release~16 also introduced \emph{conditional handover}: the network prepares one or more target
cells in advance and gives the terminal the condition; the terminal executes the handover itself when the condition is
met, removing the risk that the command arrives too late over a failing link.

\subsection{Location management and paging}

An idle terminal does not report every cell change; that would waste its battery and the network's signalling. Instead,
cells are grouped into \term{location areas} (GSM; routing areas for packet service) or \term{tracking areas} (LTE, NR).
An idle terminal reports its location only when it enters an area not in its registered list, and periodically. When a
call or data arrives, the network \term{pages} the terminal in every cell of its registered area---like calling a
friend's name across every room of a house because you only know which house they are in. Larger areas mean
fewer location updates but more paging, since each page is broadcast in more cells; smaller areas mean the reverse. The
optimum balances the two signalling loads given the users' mobility and the incoming call rate
(Figure~\ref{fig:ch20_paging}), and LTE's \emph{tracking area lists}, which let each terminal be registered in several
areas centred on its last position, avoid the storms of updates that fixed boundaries produce when a train full of
passengers crosses them. NR added a third state, RRC\_INACTIVE, in which the terminal keeps its context in the radio
network and moves within a \emph{RAN notification area} without signalling, so that it can resume data transfer quickly.

\mdcfig[0.5]{ch20_paging}{The tracking-area trade-off in a simple fluid model: location updates fall as the
square root of the area (users cross its border less often), while paging grows in proportion to it. The
total has a broad minimum; the numbers are illustrative.}

A paged terminal must be listening. To save energy, idle terminals sleep and wake only at their \emph{paging occasions},
set by a discontinuous-reception (DRX) cycle, typically from 0.32 to 2.56\,s for phones. IoT devices stretch this much
further with extended DRX and power-saving mode, in which a device can be unreachable for hours, a trade of latency for
battery life measured in years.

%======================================================================
\section{Scheduling}
\label{sec:ch20:scheduling}
%======================================================================

In an OFDMA system, a \term{scheduler} at the base station decides, for every slot and every resource block, which user
transmits, and with what modulation, coding and power. In LTE the decision is made every millisecond; in NR, every slot
of 1\,ms down to 0.125\,ms or less. That is a thousand or more decisions a second, for every cell, for hundreds of users.
The scheduler is not standardised: 3GPP specifies the signalling, and vendors compete on the algorithm. It is where the
theory of this chapter becomes money.

\subsection{Three basic policies}

Let $r_k(t)$ be the rate user $k$ could get if scheduled in slot $t$ (from its reported channel quality) and $\bar R_k(t)$
its average throughput so far.
\begin{itemize}
\item \textbf{Round robin} serves users in turn, ignoring the channel. It is fair in time, but it schedules each user
whenever its turn comes, in fades as often as peaks.
\item \textbf{Maximum rate} (max C/I) serves $\arg\max_k r_k(t)$. It maximises cell throughput, but users at the cell edge,
whose rates are always lower, are starved.
\item \textbf{Proportional fair} (PF) serves
\begin{equation}
k^*(t)=\arg\max_k\frac{r_k(t)}{\bar R_k(t)},\qquad
\bar R_k(t+1)=\Big(1-\frac1{t_c}\Big)\bar R_k(t)+\frac1{t_c}\,r_k(t)\,\mathbb 1\{k=k^*\},
\label{eq:ch20:pf}
\end{equation}
where $t_c$ is the averaging window in slots. A user is scheduled when its channel is good \emph{relative to its own
average}.
\end{itemize}

\begin{analogy}[Serving each customer at their best]
A tailor has three clients: a sprinter, a jogger and a walker. ``Max rate'' measures only the sprinter, because she
gets through fittings fastest. ``Round robin'' takes them in strict rotation, even when one has just arrived out of
breath. Proportional fair calls whichever client is \emph{relatively} sharpest today compared with their own usual
form: the walker on her best day beats the sprinter on an average one. Everyone gets fitted, and everyone is fitted
at their best. The analogy's limit: radio ``good days'' last milliseconds, so the tailor must decide a thousand
times a second.
\end{analogy}

\mdcfig{ch20_pf_trace}{The schedulers at work on three users with mean SNRs of 15, 8 and 0\,dB in slowly fading
Rayleigh channels (lines: rate each user could get; dots: who was served). Max rate serves the near user almost
every slot and never the edge user. Proportional fair spreads the slots and catches each user near its own
peaks, so even the edge user is served when its channel is unusually good.}

\subsection{Why proportional fair}

Proportional fairness comes from economics. Kelly (1998) defined an allocation of throughputs $\bar R_k$ as
proportionally fair if it maximises the sum of logarithms,
\begin{equation}
U(\bar{\vect R})=\sum_k\log\bar R_k .
\end{equation}
The logarithm has diminishing returns: doubling any user's throughput adds the same utility, whether the user has
10\,kb/s or 10\,Mb/s. An equivalent characterisation is that, at the optimum, any feasible change makes the sum of the
\emph{relative} changes non-positive: $\sum_k\Delta\bar R_k/\bar R_k\le0$.

\begin{deeper}[PF as a gradient step]
A scheduler cannot optimise long-term averages directly; it makes one decision per slot. Take a gradient step. Serving
user $k$ in slot $t$ increases $\bar R_k$ by approximately $r_k(t)/t_c$, which increases $U$ by
\begin{equation}
\Delta U\approx\frac{\partial U}{\partial\bar R_k}\cdot\frac{r_k(t)}{t_c}=\frac{1}{t_c}\,\frac{r_k(t)}{\bar R_k(t)} .
\end{equation}
The greedy choice that increases the utility most is therefore~\eqref{eq:ch20:pf}. Stochastic-approximation arguments
(Kushner and Whiting, 2004) show that for large $t_c$ this greedy rule converges to the allocation that maximises
$\sum\log\bar R_k$ over the region of achievable average throughputs. The same derivation with the utility
\begin{equation}
U_\alpha=\sum_k\frac{\bar R_k^{\,1-\alpha}}{1-\alpha}
\end{equation}
gives the \term{alpha-fair} metric $r_k/\bar R_k^{\,\alpha}$: $\alpha=0$ is max rate, $\alpha=1$ is PF (the limit of
$U_\alpha$ is the logarithm), and $\alpha\to\infty$ approaches max--min fairness, which equalises throughputs.
\end{deeper}

Two consequences are worth remembering. When all users' channels are statistically identical apart from their means,
PF gives each user an equal share of the slots, but serves each one near its own peaks. And a single knob, the
fairness exponent $\alpha$, slides the scheduler continuously from greed ($\alpha=0$) through proportional fairness
($\alpha=1$) to strict equality ($\alpha\to\infty$).

\subsection{Multiuser diversity}

The gain of channel-aware scheduling has a name: \term{multiuser diversity}. Knopp and Humblet observed in 1995 that
on a fading uplink, the sum capacity is maximised by letting only the user with the best channel transmit in each
instant. With $K$ users whose channels fade independently, the best of them is very likely to be near a peak---like
a fishing fleet that always sends out whichever boat happens to be over the shoal. For Rayleigh fading with equal mean
SNR $\bar\gamma$, the strongest of $K$ channel gains $|h_k|^2$ (independent unit exponentials) has mean
\begin{equation}
\E\Big[\max_k|h_k|^2\Big]=\sum_{k=1}^{K}\frac1k\approx\ln K+0.577,
\end{equation}
so at low SNR, where rate is proportional to SNR, the scheduler gains a factor $H_K$, 2.9 (4.7\,dB) for ten users; at
high SNR the gain is about $\log_2H_K$ b/s/Hz. Figure~\ref{fig:ch20_scheduling}a shows the cell throughput growing with
the number of users while round robin stays flat. Multiuser diversity turns fading, the enemy of the single link,
into a resource. It needs three things: channels that fluctuate (a static channel offers no peaks; Viswanath, Tse and
Laroia's \emph{opportunistic beamforming} of 2002 even induces fluctuations deliberately with randomly varying
antenna weights); fresh channel-quality feedback, so that the scheduler knows where the peaks are; and traffic that can
wait a few slots for a peak.

\mdcfig{ch20_scheduling}{Channel-aware scheduling. (a) Multiuser diversity: cell throughput versus the number of users
for the max-rate scheduler (identical to PF when users are statistically identical), with i.i.d.\ Rayleigh fading at
three mean SNRs; dotted lines show round robin. (b) The throughput--fairness trade-off for ten users with mean SNRs
spread from 20 to $-5$\,dB under the $\alpha$-fair scheduler; round robin is both less efficient and less fair than PF.}

Figure~\ref{fig:ch20_scheduling}b shows the throughput--fairness trade-off for ten users with very different mean SNRs.
Max rate ($\alpha=0$) delivers about 6.6\,b/s/Hz but a Jain fairness index of only 0.22: essentially only the best user is
served. PF ($\alpha=1$) gives about 4.3\,b/s/Hz with an index of 0.74, while round robin gives only 2.6\,b/s/Hz with an index of
0.67---PF beats round robin on both counts, because round robin wastes slots on users in deep fades. Larger $\alpha$ buys
more fairness at a rapidly rising cost. (Jain's index, $(\sum_kx_k)^2/(K\sum_kx_k^2)$, equals 1 when all users get the same
throughput and $1/K$ when one user gets everything.)

\begin{worked}[One proportional-fair decision]
Two users share a cell. User A, near the base station, has an average throughput of $\bar R_A=8$\,Mb/s; user B, at the
edge, has $\bar R_B=1$\,Mb/s. In the current slot their channel-quality reports allow $r_A=10$\,Mb/s and $r_B=1.5$\,Mb/s. Whom
does each scheduler serve?

Max rate serves A, since $10>1.5$. PF compares $10/8=1.25$ with $1.5/1=1.5$ and serves B: B is currently 50\% above its
average, A only 25\%. If B's channel had been at its average ($r_B=1$), the metrics would be 1.25 and 1.0 and A would be
served. Over time, PF gives B the slots in which B's channel is relatively best, so B's throughput is higher than under
round robin, where B would be served in its fades as well as its peaks, while A loses little.
\end{worked}

\begin{tryit}[Slide from greed to fairness]
In Lab 26's \emph{Proportional fair scheduling}, set \emph{Scheduler} to Max rate and watch the edge user starve; then
choose $\alpha$-fair and drag \emph{Fairness exponent $\alpha$} from 0 to 5, tracing the curve of
Figure~\ref{fig:ch20_scheduling}b. Lower \emph{Edge user's mean SNR} to $-10$\,dB to see how much fairness costs.
Lab 11's \emph{OFDMA scheduling, live} does the same across a whole resource grid: switch \emph{Policy} between Round
robin, Max C/I and Proportional fair with \emph{Users walk around} enabled, and its companion \emph{Link adaptation and
HARQ} shows how each scheduled rate is chosen and protected.
\end{tryit}

\subsection{Quality of service}

Real traffic is not an infinite backlog. Voice packets arrive every 20\,ms and are useless if late; video needs a steady
rate; web pages are bursty; IoT reports are tiny. LTE and NR attach each flow to a \emph{bearer} or \emph{QoS flow} with a
standardised class that tells the scheduler its priority, its packet delay budget and its tolerable packet error rate,
and whether it has a guaranteed bit rate (GBR). Table~\ref{tab:ch20:5qi} lists a few of NR's standardised 5G QoS
identifiers (5QI). Schedulers combine channel awareness with these requirements, for example with the
\emph{modified largest weighted delay first} rule, which multiplies the PF metric by the head-of-line delay of the
user's queue weighted by its urgency, so that a packet approaching its deadline gains priority---the hospital triage
nurse who sees the patient whose condition is worsening before the one who arrived first.

\begin{table}[htbp]\centering\small
\caption{A selection of standardised 5QI values (from 3GPP TS~23.501, Table 5.7.4-1; consult the current release for the
full table). Lower priority level means higher priority.}
\label{tab:ch20:5qi}
\begin{tabular}{@{}c l c c c l@{}}
\toprule
5QI & Resource type & Priority & Delay budget & Packet error rate & Example service\\
\midrule
1 & GBR & 20 & 100\,ms & $10^{-2}$ & conversational voice\\
2 & GBR & 40 & 150\,ms & $10^{-3}$ & conversational video\\
5 & non-GBR & 10 & 100\,ms & $10^{-6}$ & IMS signalling\\
9 & non-GBR & 90 & 300\,ms & $10^{-6}$ & buffered video, TCP (default)\\
82 & delay-critical GBR & 19 & 10\,ms & $10^{-4}$ & discrete automation\\
\bottomrule
\end{tabular}
\end{table}

\begin{pitfall}[Scheduler gains measured on the wrong traffic]
Most published comparisons of schedulers assume ``full buffer'' users who always have data. With real, bursty traffic,
many users have nothing to send at any instant, the number of users competing for each slot is small, and multiuser
diversity gains shrink accordingly; a user who arrives with one web page may finish before any averaging window
converges. Channel-quality reports also age: at 3.5\,GHz a terminal moving at 60\,km/h decorrelates its fast fading in a
few milliseconds, so the ``peak'' that was reported may be gone when the data are sent. Evaluate schedulers with
realistic traffic models, report delays, and include the feedback delay.
\end{pitfall}

%======================================================================
\section{The radio access network}
\label{sec:ch20:ran}
%======================================================================

The functions of this chapter---scheduling, power control, handover, paging, random access---have to live somewhere.
Where they live has changed with every generation, and the changes reflect the evolving balance between radio
performance, latency and the cost of computing and transport (Figure~\ref{fig:ch20_ran}). Chapter~\ref{ch:cellular}
tells the history of the generations; here is the architecture in brief.

\begin{figure}[htbp]\centering
\begin{tikzpicture}[font=\sffamily\scriptsize,
  box/.style={draw=navy,fill=navy!7,thick,rounded corners=1pt,minimum height=0.62cm,minimum width=1.3cm,align=center},
  core/.style={draw=accent,fill=accent!7,thick,rounded corners=1pt,minimum height=0.62cm,minimum width=1.6cm,align=center},
  lnk/.style={thick,navy!70},
  lab/.style={font=\sffamily\tiny,text=gray,above}]
\node[anchor=west,font=\sffamily\small\bfseries,text=navy] at (-1.2,0) {2G};
\node[box] (b1) at (0.6,0) {BTS}; \node[box] (b2) at (3.6,0) {BSC}; \node[core] (b3) at (6.8,0) {MSC\\(circuit core)};
\draw[lnk] (b1) -- node[lab]{Abis} (b2); \draw[lnk] (b2) -- node[lab]{A} (b3);
\node[anchor=west,text=gray,align=left] at (8.3,0) {BSC controls radio resources\\and handover for many BTSs};
\node[anchor=west,font=\sffamily\small\bfseries,text=navy] at (-1.2,-1.2) {3G};
\node[box] (n1) at (0.6,-1.2) {NodeB}; \node[box] (n2) at (3.6,-1.2) {RNC}; \node[core] (n3) at (6.8,-1.2) {CS + PS core};
\draw[lnk] (n1) -- node[lab]{Iub} (n2); \draw[lnk] (n2) -- node[lab]{Iu} (n3);
\node[anchor=west,text=gray,align=left] at (8.3,-1.2) {RNC: soft-handover combining,\\outer-loop power control};
\node[anchor=west,font=\sffamily\small\bfseries,text=navy] at (-1.2,-2.4) {4G};
\node[box] (e1) at (0.6,-2.4) {eNB}; \node[box] (e2) at (2.4,-2.4) {eNB}; \node[core] (e3) at (6.8,-2.4) {EPC\\(MME, S/P-GW)};
\draw[lnk] (e1) -- node[lab]{X2} (e2); \draw[lnk] (e2) -- node[lab]{S1} (e3);
\node[anchor=west,text=gray,align=left] at (8.3,-2.4) {flat: all radio functions,\\incl.\ scheduling, in the eNB};
\node[anchor=west,font=\sffamily\small\bfseries,text=navy] at (-1.2,-3.75) {5G};
\node[box] (g1) at (0.5,-3.75) {RU\\(low PHY,\\RF)}; \node[box] (g2) at (2.75,-3.75) {DU\\(high PHY,\\MAC, RLC)};
\node[box] (g3) at (4.95,-3.75) {CU\\(PDCP, RRC,\\SDAP)}; \node[core] (g4) at (7.2,-3.75) {5GC\\(AMF, UPF)};
\draw[lnk] (g1) -- node[lab]{fronthaul} (g2); \draw[lnk] (g2) -- node[lab]{F1} (g3); \draw[lnk] (g3) -- node[lab]{NG} (g4);
\node[draw=practc,fill=practc!8,thick,rounded corners=1pt,align=center,minimum height=0.55cm,minimum width=3.4cm] (ric) at (3.85,-5.25) {O-RAN near-real-time RIC\\(xApps, 10\,ms--1\,s loops)};
\draw[practc,thick,dashed] (g2.south) -- node[left,font=\sffamily\tiny,text=practc]{E2} (g2.south |- ric.north);
\draw[practc,thick,dashed] (g3.south) -- node[right,font=\sffamily\tiny,text=practc]{E2} (g3.south |- ric.north);
\node[anchor=west,text=gray,align=left] at (8.3,-3.75) {gNB split into CU, DU and RU;\\O-RAN opens the interfaces\\(split 7.2x fronthaul, E2, A1, O1)};
\end{tikzpicture}
\caption{Radio access network architectures. GSM and UMTS concentrated radio control in a controller (BSC, RNC); LTE moved
every radio function into a flat eNodeB connected to its neighbours by X2; 5G NR splits the gNodeB into a central unit, a
distributed unit and a radio unit, and the O-RAN Alliance adds open interfaces and a RAN intelligent controller that hosts
third-party optimisation applications.}
\label{fig:ch20_ran}
\end{figure}

\subsection{From controllers to flat base stations}

In GSM, the base transceiver station (BTS) was little more than radios and baseband; a \emph{base station controller}
(BSC) managed channels, power control and handover for tens to hundreds of BTSs, and the mobile switching centre (MSC)
switched circuits. UMTS kept the hierarchy: the NodeB did the physical layer and the \emph{radio network controller}
(RNC) did almost everything else, including combining the soft-handover branches of a CDMA call, which is why the RNC had
to be central. As data rates rose and schedulers needed millisecond reaction times, the controller became a bottleneck.
HSPA moved scheduling and retransmission into the NodeB, and LTE took the final step: the \emph{eNodeB} contains the whole
radio protocol stack, neighbouring eNodeBs coordinate handover and interference directly over the X2 interface, and the
core network (the evolved packet core) handles only mobility anchoring, sessions and user-plane routing. It is the shift
from a company where every decision goes to head office to one where branch managers decide on the spot.

\subsection{Splitting the gNodeB}

5G reversed part of the flattening, for a different reason: the cost of computing and the value of pooling. 3GPP studied
eight candidate points at which the protocol stack could be split between central and distributed equipment (TR~38.801)
and standardised one: the \emph{central unit} (CU) runs the latency-tolerant upper layers (PDCP and RRC, and SDAP for user
data), and the \emph{distributed unit} (DU) runs RLC, MAC and the upper physical layer, connected by the F1 interface. The CU
itself can be split into control and user planes. Below the DU, a \emph{radio unit} (RU) at the antenna does the lower
physical layer and RF. The fronthaul between DU and RU carries frequency-domain samples in the ``7.2x'' split specified by
the O-RAN Alliance over eCPRI, which needs far less capacity than the older CPRI interface, which carried time-domain
samples: CPRI needs about 1.2\,Gb/s per antenna for a single 20\,MHz LTE carrier, so a 64-antenna, 100\,MHz NR radio with
time-domain fronthaul would need hundreds of gigabits per second.

The \emph{O-RAN Alliance}, formed in 2018 by operators, specifies open interfaces between these units so that radios,
DUs and CUs from different vendors can interoperate, and adds a \emph{RAN intelligent controller} (RIC). The near-real-time
RIC runs third-party applications (``xApps'') that tune scheduling, handover and interference management on time scales
from about 10\,ms to 1\,s over the E2 interface; the non-real-time RIC, in the management layer, trains models and sets
policies on longer time scales. Combined with running the CU and DU software on general-purpose servers (\emph{virtualised}
or \emph{cloud} RAN), this turns many of the algorithms of this chapter---power control parameters, FFR plans, handover
offsets, scheduler weights---into software that an operator can change, and into a natural home for the machine-learning
methods discussed in Chapter~\ref{ch:sixg}.

\mdcfig{ch20_ran_timescales}{The decisions of this chapter sorted by how fast they must be made. Anything inside a
HARQ round trip lives in the DU next to the radios; handover and RRC decisions tolerate the CU's distance; the
near-real-time RIC tunes on 10\,ms to 1\,s; policies and planning run on hours to years.}

\chTwentyPhotoH[0.3]{ch20_tower_climb}{A guyed mast near Rochester, New York, with a ring of cellular panel antennas
part-way up and an equipment shelter at its foot. In a modern site the radio units sit beside the panels and the
baseband may be in the shelter or many kilometres away in a pooled data centre.}

\begin{inpractice}[Where the milliseconds go]
The functional split is constrained by latency. HARQ in LTE requires the acknowledgement of a transport block 4\,ms after
it is sent, so the MAC scheduler and the physical layer must be within about a millisecond of round-trip transport of each
other; this is why the MAC sits in the DU near the radios rather than in a distant data centre. The CU, which handles RRC
signalling and packet ordering with tolerances of tens of milliseconds, can be pooled in a regional data centre that serves
hundreds of cell sites, which is where the economies of scale come from. The near-real-time RIC's 10\,ms lower bound is not
arbitrary either: it is roughly the fastest loop that can be closed through a standard interface to a separate platform
without interfering with the slot-by-slot decisions that remain inside the DU (Figure~\ref{fig:ch20_ran_timescales}).
\end{inpractice}

%======================================================================
\section{Summary}
%======================================================================

\begin{itemize}
\item Sharing a band requires duplexing and multiple access. FDD offers continuity; TDD offers flexibility, unpaired
spectrum and channel reciprocity but needs a guard period of at least the round-trip delay plus switching time and
network-wide synchronisation. Full duplex needs about 110\,dB of self-interference cancellation; subband full duplex at
the base station is the practical first step.
\item FDMA, TDMA, CDMA, OFDMA and SDMA all exploit orthogonality in some dimension and are equivalent in an ideal single
cell; they differ in how they lose orthogonality and how they cope with inter-cell interference. NOMA superimposes users
with different channel strengths and achieves the broadcast-channel capacity region with SIC.
\item Pure and slotted ALOHA reach $1/(2e)$ and $1/e$, and are unstable without adaptive retransmission. Carrier sense helps
when $a=\tau/T_p$ is small; CSMA/CD reaches $1/(1+3.44a)$. On radio channels, hidden terminals defeat carrier sense;
802.11 uses CSMA/CA with binary exponential backoff and optional RTS/CTS, accurately predicted by Bianchi's fixed point.
\item Cellular random access is slotted ALOHA with $M$ Zadoff--Chu preambles per occasion: at most about $M/e$ successes per
occasion. Massive IoT bursts need barring, more resources and fewer messages.
\item Hexagonal reuse uses clusters of $N=i^2+ij+j^2$ cells, with $D/R=\sqrt{3N}$ and first-tier
$\mathrm{SIR}\approx(3N)^{n/2}/6$. AMPS needed $N=7$; sectorization and cell splitting trade cost for SIR and capacity.
\item Erlang B, $B(A,C)$, dimensions blocked-calls-cleared systems and is insensitive to the holding-time distribution;
Erlang C gives waiting probabilities. Large trunks are efficient, so splitting pools costs capacity.
\item Reuse-1 networks manage interference with fractional power control, FFR/SFR, eICIC (range expansion with almost blank
subframes) and CoMP.
\item With PPP base stations and Rayleigh fading, the interference-limited coverage is $1/(1+\rho(T,\alpha))$, independent of
density, so area spectral efficiency grows linearly with density until multi-slope path loss, antenna heights and cost
intervene.
\item Handover trades ping-pong against delay through hysteresis and time-to-trigger; idle terminals are tracked by
location or tracking areas and paged.
\item Proportional fair scheduling maximises $\sum\log\bar R_k$ by serving $\arg\max r_k/\bar R_k$; it captures multiuser
diversity (a gain of about $\ln K$ in SNR) while remaining fair. The RAN has evolved from controller-based to flat to split
and open architectures.
\end{itemize}

\begin{labbox}[Multiple access and cellular systems]
Two interactive labs accompany this chapter:
\begin{itemize}
\item \lab{moderncomms-labs/labs/lab26\_cellular\_system.py} --- \emph{Clusters, reuse and SIR} (SIR maps and CDFs for
any cluster size, with sectors and shadowing); \emph{Erlang: a live switchboard} (animated call arrivals, blocking and
the trunking penalty of sectors); \emph{ALOHA and CSMA, live} (packet timelines and throughput curves);
\emph{Slotted ALOHA's instability} (fixed versus pseudo-Bayesian retransmission); \emph{Wi-Fi's DCF: Bianchi's model}
(fixed point versus simulation, basic and RTS/CTS); \emph{RACH collisions and barring} (one occasion, massive bursts and
access-class barring); \emph{Stochastic geometry: PPP coverage} (random deployments against the closed form);
\emph{Proportional fair scheduling} (round robin, max rate and $\alpha$-fair, with Jain's index); and \emph{Handover:
hysteresis and TTT} (ping-pong against late handover).
\item \lab{moderncomms-labs/labs/lab11\_air\_interfaces.py} --- \emph{Frequency reuse and SIR} (reuse plans and the
channels each costs); \emph{Trunking: a live switchboard}; \emph{OFDMA scheduling, live} (walking users and a scheduled
resource grid); \emph{Link adaptation and HARQ}; \emph{Wi-Fi contention, live}; \emph{4096-QAM: the EVM budget}; and
\emph{Air interfaces side by side} (GSM to NR compared). The last three also serve Chapters~\ref{ch:ofdm},
\ref{ch:cellular} and~\ref{ch:wifi}.
\end{itemize}
The functions the labs use (Erlang formulas, hexagonal geometry, Bianchi's fixed point, PPP coverage, schedulers) are in
\lab{moderncomms-labs/commlib/cellular.py}, and \lab{book/figscripts/ch20\_figs.py} generates every data figure in this
chapter.
\end{labbox}

\problems
\begin{enumerate}[label=\thechapter.\arabic*]
\item An FDD band has a 20\,MHz uplink at 1920--1940\,MHz and a downlink at 2110--2130\,MHz. What are the duplex spacing and the
duplex gap? A TDD operator with a 40\,MHz unpaired block uses the pattern DDDSU with the special slot split 10:2:2 (30\,kHz
subcarrier spacing). Compare the downlink and uplink capacity of the two operators, assuming equal spectral efficiency.
\item Using~\eqref{eq:ch20:gp}, find the guard period needed for a TDD cell of radius 15\,km with a 15\,$\mu$s switching time.
How many 30\,kHz OFDM symbols is that? What fraction of a 2.5\,ms DDDSU period is lost?
\item Derive the throughput of pure ALOHA,~\eqref{eq:ch20:pure}, and of slotted ALOHA. For slotted ALOHA at $G=0.5$, 1 and 2,
compute the fractions of idle, successful and collided slots.
\item In slotted ALOHA with $m$ users each transmitting with probability $p$ in every slot, show that the success probability is
$mp(1-p)^{m-1}$, that it is maximised by $p=1/m$, and that the maximum tends to $1/e$ as $m\to\infty$. How large must $m$ be for
the maximum to be within 1\% of $1/e$?
\item Derive the nonpersistent CSMA throughput~\eqref{eq:ch20:csma} by the renewal argument of the text. Find the maximum
throughput for $a=0.001$, 0.01 and 0.1 numerically, and the offered load at which it occurs.
\item A 100\,Mb/s half-duplex Ethernet segment has a slot time of 512 bit times. What is the maximum network diameter if the
signal propagates at $2\times10^8$\,m/s and repeater and transceiver delays consume half the budget? Using~\eqref{eq:ch20:csmacd},
find the efficiency with 64-byte and 1518-byte frames.
\item Solve Bianchi's equations~\eqref{eq:ch20:bianchitau}--\eqref{eq:ch20:bianchip} by hand (or with a few iterations on a
calculator) for $n=5$, $W=16$, $m=6$. Then use~\eqref{eq:ch20:bianchiS} with $\sigma=9\,\mu$s, $T_s=330\,\mu$s,
$T_c=320\,\mu$s and $T_{\text{pl}}=222\,\mu$s to find the throughput as a fraction of the PHY rate.
\item A cell has 54 contention preambles and a RACH occasion every 10\,ms. Using~\eqref{eq:ch20:rach}, what is the largest
sustainable access rate? If 2000 devices try to access within 1\,s, uniformly, what fraction collide on their first attempt?
\item Show that the cluster sizes allowed by hexagonal geometry are $N=i^2+ij+j^2$ and list those below 30. For a path-loss exponent
of 3, find the smallest $N$ that gives a first-tier SIR of 15\,dB with omnidirectional cells, and with three $120^\circ$ sectors
(two interferers).
\item A system has 33\,MHz of spectrum with 25\,kHz duplex channel pairs (both directions), of which 21 per cluster are control
channels. Find the number of voice channels per cell for $N=4$, 7 and 12. With 2\% blocking and 30\,mE per subscriber, how many
subscribers does a cell support in each case?
\item Derive the Erlang B formula~\eqref{eq:ch20:erlangb} from the birth--death balance equations and verify the
recursion~\eqref{eq:ch20:erlangbrec}. Then derive Erlang C~\eqref{eq:ch20:erlangc} and show that it can be written in terms of
Erlang B as given.
\item Using~\eqref{eq:ch20:pppnoisefree}, compute the coverage probability of a PPP network with $\alpha=4$ at $T=-3$, 0, 3 and
10\,dB. What fraction of users have SINR above 0\,dB if the network uses reuse 3 (so that interferers form a thinned PPP of density
$\lambda/3$)? \emph{Hint:} only the interference term changes.
\item Show that the PF rule~\eqref{eq:ch20:pf} is the gradient step for $U=\sum_k\log\bar R_k$, and that the $\alpha$-fair rule is
the gradient step for $U_\alpha$. For two users with i.i.d.\ Rayleigh channels and mean SNRs of 20 and 0\,dB, which user gets more
slots under PF? Under max rate?
\item \emph{(Simulation)} Write an event-driven simulation of slotted ALOHA with 100 users, Bernoulli arrivals and (a) fixed
retransmission probability, (b) binary exponential backoff, (c) Rivest's pseudo-Bayesian rule. Plot throughput and mean delay
against the arrival rate and find where each becomes unstable.
\item \emph{(Simulation)} Reproduce Figure~\ref{fig:ch20_ppp}b with a hexagonal grid and with a PPP, then add log-normal shadowing
with $\sigma=8$\,dB and repeat. Show that shadowing makes the hexagonal network's SINR distribution approach the PPP result, and
explain why.
\end{enumerate}

\furtherreading
\begin{itemize}
\item V.~H.~MacDonald, ``Advanced Mobile Phone Service: The cellular concept,'' \emph{Bell System Technical Journal}, vol.~58,
no.~1, pp.~15--41, 1979. The classic exposition of hexagonal reuse, cluster sizes and cell splitting, in the issue that described
AMPS.
\item T.~S.~Rappaport, \emph{Wireless Communications: Principles and Practice}, 2nd ed., Prentice Hall, 2002. Chapter~3 on the
cellular concept, interference and trunking, with worked Erlang examples, is the standard undergraduate treatment.
\item D.~Bertsekas and R.~Gallager, \emph{Data Networks}, 2nd ed., Prentice Hall, 1992. Chapter~4 on multiaccess communication
covers ALOHA, its stability, pseudo-Bayesian control, splitting algorithms, CSMA and CSMA/CD with exemplary rigour.
\item N.~Abramson, ``The ALOHA system---another alternative for computer communications,'' \emph{Proc.\ AFIPS Fall Joint
Computer Conference}, 1970; and R.~M.~Metcalfe and D.~R.~Boggs, ``Ethernet: Distributed packet switching for local computer
networks,'' \emph{Communications of the ACM}, vol.~19, no.~7, 1976. Two founding papers, both short and readable.
\item L.~Kleinrock and F.~A.~Tobagi, ``Packet switching in radio channels: Part I---Carrier sense multiple-access modes and
their throughput-delay characteristics,'' \emph{IEEE Transactions on Communications}, vol.~23, no.~12, 1975. The CSMA analysis; Part~II
(Tobagi and Kleinrock, same journal and year) introduced the hidden-terminal problem.
\item G.~Bianchi, ``Performance analysis of the IEEE 802.11 distributed coordination function,'' \emph{IEEE Journal on Selected
Areas in Communications}, vol.~18, no.~3, 2000. The fixed-point model of Section~\ref{sec:ch20:random}.
\item J.~G.~Andrews, F.~Baccelli and R.~K.~Ganti, ``A tractable approach to coverage and rate in cellular networks,'' \emph{IEEE
Transactions on Communications}, vol.~59, no.~11, 2011; and M.~Haenggi, \emph{Stochastic Geometry for Wireless Networks},
Cambridge University Press, 2012. The PPP model of cellular networks and the mathematics behind it.
\item D.~Tse and P.~Viswanath, \emph{Fundamentals of Wireless Communication}, Cambridge University Press, 2005. Chapter~4 on
cellular systems and Chapter~6 on multiuser capacity, opportunistic communication and proportional fair scheduling.
\item F.~P.~Kelly, A.~K.~Maulloo and D.~K.~H.~Tan, ``Rate control for communication networks: shadow prices, proportional fairness
and stability,'' \emph{Journal of the Operational Research Society}, vol.~49, 1998; and J.~Mo and J.~Walrand, ``Fair end-to-end
window-based congestion control,'' \emph{IEEE/ACM Transactions on Networking}, vol.~8, no.~5, 2000, which introduced
$\alpha$-fairness.
\item E.~Dahlman, S.~Parkvall and J.~Sk\"old, \emph{5G NR: The Next Generation Wireless Access Technology}, 2nd ed., Academic Press,
2020. Random access, scheduling, power control, mobility and the RAN architecture of NR, by members of the team that designed it.
\item L.~Kleinrock, \emph{Queueing Systems, Volume 1: Theory}, Wiley, 1975. Birth--death processes, Erlang's formulas and the
M/M/$C$ queue, from first principles.
\end{itemize}

